% =====================================================================
%  Capitolo 15 — La revisione di compendium e vision di Metroid Dread
%  (pilastri, must do, must not, nice to have) e le decisioni nel gioco:
%  tipi di decisione e le tre trappole (meaningless, obvious, blind)
%  (Game Design, primo anno)
%  Fonte: trascrizione della lezione GD15 (Game Design, primo anno),
%  stesso docente dei capitoli 9 e 10 (corso di design methods), di cui
%  la lezione riprende il filo dopo i quattro capitoli di level design.
%  La lezione si è svolta da remoto, su Zoom, con una lavagna condivisa
%  su Miro.
%  Contenuti, nell'ordine della lezione: l'apertura da remoto e la
%  lavagna Miro; il programma (design methods, una lezione importante
%  venerdì prossimo, il deep dive sui flowchart); la comunicazione
%  ufficiale sulle intrusioni esterne e l'invito a parlare; la webcam;
%  la revisione dei task su Metroid Dread: dove si documentano le
%  posizioni di nemici e oggetti (domanda di Valerio: il level design
%  document, la mappa 2D e il blockout, il Master Asset Document, lo
%  hero asset); lo stealth non è un pilastro, gli EMMI sì (le zone
%  EMMI come spina dorsale, lo stealth come strategia e non meccanica,
%  gli unique selling point); il target audience e il doppio senso di
%  ``estetiche''; i must do, must not e nice to have lato design (la
%  colpa del mock-up, il counter, la camera, le abilità, gli EMMI, i
%  biomi, la mappa, le puntine di Hollow Knight, l'economy sink); il
%  tempo dedicato alla ricerca (Federica); dividere i must per
%  dipartimento e la vestizione (Francesca; le tre camere di Zelda, il
%  numero di abilità e i costi, lo skill tree); gli esempi non sono la
%  regola; il task dei feedback su Zoe; l'annuncio del flowchart di
%  Grotte e i due errori più comuni; il calendario e il task di
%  pitching; la pausa; il Beer to Peer; le tre porte; che cos'è una
%  decisione; i tipi di decisione (tema, meccaniche, meta); le tre
%  trappole (meaningless, obvious, blind) con BioShock Infinite,
%  l'Anfora dell'Avidità e i doni iniziali di Dark Souls; la
%  discussione sulla ricompensa emotiva (Federica, Valerio, Francesca:
%  Silksong, Life is Strange, The Walking Dead di Telltale); le
%  decisioni ben fatte (Oxenfree, The Witcher, Fire Emblem, Pokémon,
%  Life is Strange); il riepilogo delle leve per tipo; il task (un
%  gioco breve dello stesso autore di The Stanley Parable) e la
%  chiusura.
%  Integrazioni da "Challenges for Game Designers" (Brathwaite e
%  Schreiber), "The Art of Game Design" (Jesse Schell) e "An
%  Architectural Approach to Level Design" (Christopher W. Totten),
%  segnalate come approfondimenti.
%
%  Pacchetti richiesti nel preambolo del libro (già presenti nel
%  main.tex): \usepackage{booktabs}, \usepackage{amsmath},
%  \usepackage{enumitem}, \usepackage{tikz} con
%    \usetikzlibrary{shapes.geometric,shapes.symbols,shapes.misc,
%      arrows.meta,positioning,calc,fit}
%  (le librerie vanno caricate nel preambolo, non nel capitolo).
%  Tutte le etichette di questo capitolo hanno il prefisso "gd15:"
%  (gd15:sec:..., gd15:fig:..., gd15:tab:...) e tutti gli stili TikZ
%  il prefisso "gd15", per non collidere con gli altri capitoli.
%
%  Nota sulla fonte: la trascrizione automatica è buona nelle parti
%  discorsive, ma storpia moltissimi nomi propri e termini inglesi
%  (``Mero'' per Miro, ``MassDu''/``Mass2'' per must do, ``Nice2EV''
%  per nice to have, ``Unix all-end point'' per unique selling point,
%  ``Enni''/``M''/``Lemmy'' per gli EMMI, ``Holonite''/``Ullonite'' per
%  Hollow Knight, ``Singsong''/``Salesong''/``Silzong'' per Silksong,
%  ``Kazin'' per cutscene, ``Aux Free'' per Oxenfree e così via). Dove
%  il riferimento è ricostruibile con certezza lo si riporta nella
%  forma corretta; dove non lo è, lo si segnala in nota invece di
%  indovinare. Come nei capitoli 9 e 10, durante la pausa la
%  registrazione ha trascritto in loop una frase senza senso
%  (``Grazie a tutti per l'attenzione... Qual è la risposta che ci può
%  essere?''): è un artefatto, omesso. Le slide non sono nella fonte:
%  figure e tabelle sono ricostruzioni schematiche di ciò che il
%  docente descrive e mostra, non le slide originali.
% =====================================================================

\chapter[Vision, must have e decisioni nel gioco]{La revisione della vision
e le decisioni nel gioco: meaningless, obvious, blind}
\label{gd15:cap:decisioni}

Dopo i quattro capitoli di level design (dal capitolo~\ref{gd8:cap:leveldesign}
in avanti), si torna al corso di \emph{design methods} con lo stesso docente
dei capitoli~\ref{gd9:cap:agency-visione} e~\ref{gd10:cap:feedback}. La
lezione si svolge \textbf{da remoto}, su Zoom, per ragioni esterne (``avrei
gradito fare lezione in presenza, come potete immaginare; purtroppo non mi è
stato possibile''), e ha due anime ben distinte, che il capitolo rispecchia:
\begin{enumerate}
  \item una lunga \textbf{prima parte di revisione}: le comunicazioni di
        servizio, poi i dubbi sui task di \textbf{compendium e vision di
        \emph{Metroid Dread}} (capitolo~\ref{gd9:cap:agency-visione}),
        lasciati apposta aperti nei commenti per discuterne a lezione: dove
        si documentano le posizioni di nemici e oggetti, perché lo stealth
        non è un pilastro ma gli EMMI sì, il target audience, e soprattutto
        che cosa va nei \textbf{must do}, nei \textbf{must not} e nei
        \textbf{nice to have} dal punto di vista del design; infine il task
        dei feedback su ``Zoe'' e l'annuncio di un flowchart da fare insieme
        (sezioni~\ref{gd15:sec:apertura}--\ref{gd15:sec:calendario});
  \item dopo la pausa, la \textbf{lezione nuova}: le \textbf{decisioni}.
        Che cos'è una decisione in un gioco, quali \textbf{tipi} di
        decisione esistono (legate al tema, alle meccaniche, meta), quali
        sono le \textbf{tre trappole} da evitare (decisioni
        \emph{meaningless}, \emph{obvious} e \emph{blind}), come si fanno
        invece le decisioni ben costruite e con quali leve; poi il task
        (sezioni~\ref{gd15:sec:porte}--\ref{gd15:sec:task}).
\end{enumerate}
Il flowchart collettivo, previsto per l'ultima mezz'ora, alla fine
\textbf{non si farà}: la discussione sulle decisioni si prende tutto il
tempo, e il docente lo rimanda alla lezione successiva
(sezione~\ref{gd15:sec:prossima}).

% ---------------------------------------------------------------------
\section{In apertura: una lezione da remoto}
\label{gd15:sec:apertura}

\subsection{Presenze e una lavagna condivisa}

Il docente saluta gli studenti man mano che si collegano (Mattia, Francesca,
Alberto, Federica, poi Valerio e, in ritardo, Daniele), ammette di non aver
segnato le presenze la volta precedente e chiede chi manca: di sicuro
Lorenzo, forse Francesco. In attesa degli altri chiede a tutti di entrare in
un \textbf{link} a una lavagna condivisa su \textbf{Miro}:\footnote{Nella
trascrizione ``Mero''; dal contesto (una board condivisa in cui ``scrivere e
fare disegni'', che chiede la creazione di un account e un'email di lavoro)
si tratta quasi certamente di Miro, la lavagna collaborativa online.} può
chiedere di creare un account, e in effetti a qualcuno risponde che il
progetto è privato. Il docente cambia le impostazioni di condivisione (``è
uno dei tanti motivi per cui io odio profondamente'' questo strumento: per
l'accesso chiede un'email di lavoro), e alla fine tutti riescono a entrare,
scrivere e disegnare. Valerio aveva chiesto l'accesso via email, ma è
riuscito a entrare comunque; a Daniele, arrivato dopo e che non vede la
cronologia della chat di Zoom, il docente rimanda i messaggi.\footnote{Nella
trascrizione, durante l'attesa, compare anche un frammento senza alcun
legame con la lezione (``la situazione è questa, la migliore dal 2005
all'inizio della storia della città\ldots la guerra civile\ldots''): è con
ogni probabilità audio esterno o un errore del riconoscimento vocale, e viene
omesso.}

\subsection{Il programma della giornata}

Prima ancora di condividere lo schermo, il docente ``smarca'' le cose
veloci.
\begin{itemize}
  \item Oggi si fa ancora \textbf{design methods}, non
        \emph{documentation and pitching}.\footnote{Nella trascrizione ``non
        documentation and peaching, ma design methods''. ``Peaching'' è
        quasi certamente \emph{pitching}: le due anime del corso del docente
        sono la documentazione e la presentazione da un lato e i metodi di
        design dall'altro.}
  \item Dopo questa lezione si entra in un nuovo blocco, che parte con la
        lezione di \textbf{venerdì prossimo}: una lezione
        \textbf{particolarmente importante}, che il docente chiede di seguire
        in presenza ``se possibile'', anche perché vuole discutere di
        altre questioni (sezione~\ref{gd15:sec:comunicazioni}).
  \item Avendo notato \textbf{una serie di problemi sui flowchart} (lo aveva
        scritto anche su Discord), il docente vuole fare un \textbf{deep
        dive} e un \textbf{flowchart tutti insieme}, come era stato fatto
        all'inizio del corso con il documento di formattazione, su un gioco
        più contenuto e semplice. ``Non ho mai nascosto'' che il task di
        tradurre \emph{Pac-Man} in flowchart (sezione~\ref{sec:task-gd6})
        ``non era semplice''.
  \item Lo stesso vale per il task su \emph{Dread}: non tanto il compendium
        quanto la \textbf{vision}, su cui il docente ha lasciato commenti
        apposta per aprire discussioni in aula. Il compendium e la vision si
        discutono \textbf{per primi}, il flowchart dopo.
\end{itemize}

% ---------------------------------------------------------------------
\section{Comunicazioni: niente più intrusioni, e parliamone}
\label{gd15:sec:comunicazioni}

\subsection{L'informazione ufficiale}

C'è una sola informazione ufficiale (per il resto il docente rimanda alla
lezione di venerdì, ``la prossima volta che mi ritroverete in carne ed
ossa''), e dovrebbe essere già arrivata anche con un comunicato: gli studenti
\textbf{non avranno più intrusioni} in nessuna lezione del primo anno.
Che ci sia lui, Alessandro, Stefano, oppure Daniele o Giacomo per le lezioni
di narrativa, nessun membro esterno entrerà più in aula, \textbf{a meno che
non lo richieda l'intera classe}: basta che una sola persona non sia
d'accordo, e serve il lasciapassare di tutti. È una cosa che gli studenti
avevano chiesto, che è passata dal corpo docenti e che è stata risolta
``anche fin troppo velocemente''; il docente si dice d'accordo con
``tutto ciò che è positivo per voi e che voi richiedete''.

\subsection{Un invito a parlare}

Il docente ha parlato velocemente con Alessandro (Alessandro Pregnolato, il
docente di programmazione, ``che magari non era chiaro'') di alcune
questioni emerse, quindi non ha informazioni ``a 360 gradi''. Ma alcune di
queste problematiche lui è ben disposto a discuterle, perché la scuola cerca
sempre di venire incontro alle esigenze degli studenti, soprattutto per il
loro benessere. Per questo li invita a \textbf{parlarne}, anche con lui e
non solo con altri docenti: da remoto è più difficile, quindi rimanda a
venerdì, quando sarà presente sia durante la lezione sia dopo.

Una delle questioni riguarda \textbf{il rapporto con il docente durante la
lezione registrata}. La registrazione è ``uno strumento sempre per voi'':
serve a riprendere le lezioni o a far interagire chi non può essere presente
(o, come oggi, a fare lezione quando non si può farla in presenza per
fattori esterni). Gli studenti non devono preoccuparsi di visualizzazioni
esterne, che sarebbero illegali, né di osservazioni da parte di altri
docenti: ``non funziona così''. Il docente è sempre disponibile a qualsiasi
discorso, dentro o fuori lezione; gli studenti possono anche dirgli di no,
``però vi chiederei di dirmi di no'', cioè di dirlo esplicitamente.

\subsection{Tenete accesa la webcam}

Sempre per via della lezione da remoto, il docente chiede a tutti di
\textbf{tenere accesa la webcam}: gli serve per capire se gli studenti lo
stanno seguendo, ``se siete vivi'', per vedere le facce, le perplessità,
``qualche ciglio alzato''. La richiesta vale per tutti i docenti, non solo
per lui: da remoto è molto difficile capire se si è capito ``o se vi state
addormentando''. Se qualcuno non ha la webcam, pazienza: sono casi limite.

% ---------------------------------------------------------------------
\section{Il compendium: dove stanno le posizioni di nemici e oggetti?}
\label{gd15:sec:posizioni}

\subsection{La domanda di Valerio}

Il docente ha commentato tutto, anche i feedback sugli \emph{zombie game}
(sezione~\ref{gd15:sec:zoe}), e lascia la parola agli studenti. Valerio
apre: nel compendium aveva previsto documenti per le \textbf{location dei
nemici e degli oggetti}, e il commento diceva che era ``un caso in cui le
cose descritte vanno insieme'', da chiarire a lezione. Il docente condivide
lo schermo e apre il task.

\subsection{Documenti accorpati}

Molti studenti hanno fatto la stessa cosa: documenti separati, sia nella
categoria di game design sia in quella di level design, per ogni aspetto.
La risposta del docente parte da una premessa già vista
(sezione~\ref{gd9:sec:compendium}): nelle produzioni moderne il \textbf{GDD
completo} del gioco non è più legato tanto all'operatività, cioè a
\emph{fare} il gioco, quanto a \emph{spiegarlo}: appartiene più al reame
della produzione, della visione o del generale, ``qualsiasi altra cosa
vogliate chiamarla''.
\begin{itemize}
  \item Nel \textbf{game design} è normale che esistano documenti che
        uniscono o spiegano \textbf{tante meccaniche} insieme: una board di
        Miro, oppure un documento sul \emph{combat system} generale che
        contiene molte meccaniche spiegabili singolarmente.
  \item Nel \textbf{level design} la cosa diventa più problematica. Il
        docente premette di avere meno esperienza lato level di, per
        esempio, Stefano; ma anche nelle produzioni con un level design
        ``molto massiccio'' non si parla mai di \textbf{documenti specifici}
        per le location o per gli oggetti. Si ha di solito un \textbf{level
        design document generale}, oppure quello \textbf{di un livello} o
        \textbf{di un bioma}, che contiene al suo interno tutti questi
        piazzamenti.
\end{itemize}

\subsection{Dalla mappa 2D al blockout, e la discrepanza}

Il docente chiede a Valerio se intendesse la posizione specifica (per
esempio dei punti delle spline lungo cui si muovono i nemici) o il numero
generale di nemici per livello. Valerio intendeva la \textbf{posizione
specifica}: e lì la documentazione diventa ancora più problematica.

Il level design, come il game design, è una disciplina in cui si fa
\textbf{documentazione a monte}: c'è un level design document, o una
\textbf{mappa 2D} (tra gli strumenti ce ne sono tanti, per esempio la
\textbf{bubble map}). Poi la mappa la si \textbf{riporta in engine con un
blockout}, esattamente come Stefano ha chiesto agli studenti di fare per la
loro stanza (sezione~\ref{gd11:sec:blockout}): in engine, senza
documentazione di per sé. Ma tra le due cose c'è \textbf{quasi sempre una
discrepanza molto grossa}: ciò che si era previsto sulla mappa 2D, quando si
passa al 3D (a meno che il gioco non sia top-down), va spesso modificato.

Sulla carta si avrà quindi un'idea del \textbf{numero generico} di nemici e
di \textbf{dove} dovrebbero essere: si sa che il boss è in quella stanza, o
che prima della porta ci devono essere due guardie. Ma \textbf{che percorsi
fanno} o \textbf{quanto sono lontane} l'una dall'altra è qualcosa che si
vede solo in engine, e che è molto più facile modificare lì senza dover poi
aggiornare il documento. I documenti di level design accorpano dunque tante
cose, il cui \textbf{dettaglio sta in engine}.

\paragraph{Approfondimento: i diagrammi di prossimità.} La bubble map citata
dal docente ha un parente stretto in architettura, che Totten descrive in
\emph{An Architectural Approach to Level Design}: il \textbf{diagramma di
prossimità}. Come i diagrammi a molecola, è fatto di bolle collegate da
linee: le bolle sono gli spazi richiesti (in un livello: la stanza del boss,
la postazione da cecchino, il traguardo), dimensionate in base alla loro
importanza o grandezza, e le linee dicono quanto è importante che due spazi
siano vicini e che tipo di connessione hanno. Totten insiste però su un
punto che è esattamente la ``discrepanza'' del docente: questi diagrammi
\textbf{non sono piante}, ma strumenti per analizzare l'idea funzionale; il
risultato finale ``può (e deve) avere un aspetto drasticamente diverso''.

\subsection{Quando un oggetto sta nel game design}

C'è una differenza da fare. Per il \textbf{funzionamento} di porte o di
meccaniche come gli \textbf{oggetti di vaulting}, la documentazione sta di
solito nel \textbf{game design}. L'esempio: c'è una staccionata alta tot
unità di Unreal; se il character si avvicina esce un \emph{prompt}; se il
giocatore preme il tasto parte l'animazione, il salto e l'arrivo dall'altra
parte. Il funzionamento di questo oggetto in gioco l'ha definito qualcuno
del game design, anche se oggetti del genere stanno a cavallo dei due
dipartimenti. Questa è l'unica eccezione.

\subsection{Il Master Asset Document}

Il \textbf{numero} di porte in una data zona o sezione del livello, o il
numero di oggetti di gioco (nemici inclusi), sta invece di solito in un
\textbf{documento unico}. La nomenclatura non è sempre condivisa, ma nella
sua esperienza il termine più ricorrente è \textbf{MAD}, acronimo di
\textbf{Master Asset Document}: per ogni zona specifica del livello
racchiude ``proprio una lista, un Excel enorme'' (gli studenti non lo
vedranno) con tutti gli elementi necessari e i \textbf{dipartimenti
coinvolti}. Per esempio:
\begin{itemize}
  \item c'è l'oggetto di vault, e servirà chiaramente la sua
        \textbf{animazione specifica};
  \item c'è un cosiddetto \textbf{hero asset}: un oggetto unico, modellato
        specificamente, che fa da punto guida, \textbf{punto di interesse},
        \textbf{landmark} per orientare il giocatore.
\end{itemize}
Tutte queste cose sono necessarie, ma \textbf{molto più accorpate che
separate}, e spesso più in engine che in un documento. Come sempre, conclude
il docente, l'obiettivo del task non era la precisione di un momento della
produzione che gli studenti non hanno ancora affrontato, ma iniziare a farli
ragionare. Valerio conferma che ora è più chiaro.

% ---------------------------------------------------------------------
\section{I pilastri: lo stealth no, gli EMMI sì}
\label{gd15:sec:pilastri}

\subsection{Le due cose difficili del task}

Un'altra studentessa ricorda due commenti rimasti aperti: uno sui
\textbf{pilastri}, in cui aveva indicato lo \textbf{stealth} e il docente
aveva risposto che ``discutibilmente'' era un pilastro di \emph{Dread}
(``non sono completamente d'accordo''), e uno sul \textbf{target audience}
(``apprezzo il tentativo, ma ne riparliamo''). Il docente parte dai pilastri
e confessa, in retrospettiva, che le due cose davvero difficili del task
erano i \textbf{pilastri} (sezione~\ref{gd9:sec:pilastri}), ``tosti a
prescindere'', e gli \textbf{unique selling point}.\footnote{Nella
trascrizione ``gli unix all-end point'', poi ``Unix Endpoint'' e ``unix
all'inpoint'': come nel capitolo~\ref{gd10:cap:feedback}, è una storpiatura
di \emph{unique selling point} (USP).}

\subsection{Esplorazione e combattimento}

Due pilastri sono sicuri, e molti li hanno colti, anche con formulazioni
diverse: \textbf{esplorazione} e \textbf{combattimento}.
\begin{itemize}
  \item Sul \textbf{combattimento}, Francesca è rimasta vicina all'esempio
        di \emph{God of War} visto a lezione (tabella~\ref{gd9:tab:gow}).
        Ma il combattimento di \emph{Dread} è lontano da quello di
        \emph{God of War}: è un sistema \textbf{molto più snello}, anche se
        è vero ciò che hanno scritto lei e altri sul \textbf{melee}, sul
        \textbf{counter} e sul modo più ravvicinato di combattere di Samus
        (che ``guarda caso'' sembra ancora più spinto nel prossimo
        capitolo).
  \item Sull'\textbf{esplorazione} ``niente da fare'': si parla di un
        \emph{metroidvania}, con la sua parte esplorativa e il suo mondo
        interconnesso.
\end{itemize}

\subsection{Perché lo stealth non è un pilastro}

Sullo stealth il docente non ha la risposta, ma azzarda un'ipotesi: lato
marketing può funzionare, e probabilmente \textbf{voleva essere più spinto}
di quanto sia risultato. Il gioco dà al giocatore, come secondo potenziamento
(il primo, generalmente, è la \textbf{Morph Ball},\footnote{Nella
trascrizione ``la morfosfera''. Il docente aggiunge subito che ``la
morfosfera è da mettere nei must do'': ci si torna nella
sezione~\ref{gd15:sec:must}.} il \textbf{Phantom Cloak}:\footnote{Nella
trascrizione ``cloak stealth'' e poi ``Phantom Clock''; è il Phantom Cloak,
il mantello che rende Samus invisibile.} se stai fermo, diventi invisibile.
Funziona, sì: per aprire un paio di porte e per non farsi scoprire dagli
\textbf{EMMI} nelle loro zone.\footnote{Nella trascrizione gli EMMI (i robot
cacciatori di \emph{Metroid Dread}) compaiono come ``Emmy'', ``Enni'',
``Lemmy'', ``M'' ed ``Envi''.} Ma:
\begin{itemize}
  \item quasi sempre le \textbf{zone EMMI} sono pensate in modo che il
        giocatore \textbf{entri e poi esca}, più che si nasconda;
  \item il cloak \textbf{smette di servire} a un certo punto: funziona solo
        se sai già dove si trova l'EMMI, perché quando l'EMMI è \textbf{in
        allerta}, anche se diventi invisibile ti trova comunque: scatta lo
        stato di caccia;
  \item non è così potente, e dalla metà dell'esperienza in poi viene
        percepito come ulteriormente depotenziato.
\end{itemize}

\subsection{Le zone EMMI come spina dorsale}

Il docente apre la mappa di \emph{Metroid Dread} e la percorre dall'inizio
(figura~\ref{gd15:fig:emmi}). Il giocatore parte più o meno in una certa
zona, fa un percorso con un'interazione narrativa, sale, trova il
salvataggio, poi \textbf{entra nella zona EMMI} (c'è una tutorializzazione
iniziale con gli EMMI, ma la si ignora), \textbf{esce dalla zona EMMI},
salva, sale, \textbf{rientra} nella zona EMMI, \textbf{riesce}. E questo
viene fatto $N$ volte in quasi tutti gli ambienti. È una cosa che, ricorda il
docente, aveva scritto anche Lorenzo nel suo task, sebbene in modo diverso.

Le zone EMMI sono ``come se fossero \textbf{la spina dorsale del livello}'':
sono poste sempre a metà, in modo da \textbf{tagliare in due} il livello, così
che il giocatore sia spinto a entrarvi non solo per proseguire, ma anche
perché, ``permettetemi il francese'', \textbf{ha paura}: non può uccidere
questi mostri, deve scappare verso la via di fuga più vicina appena viene
visto. Già alla terza volta, però, il giocatore capisce come funziona lato
strategia: l'EMMI non è più tanto il nemico unico della zona quanto un
problema di \textbf{uscire dalla zona il prima possibile}. Non devo evitare
l'EMMI: devo proprio scappare.

\begin{figure}[htbp]
  \centering
  \begin{tikzpicture}[
      gd15zona/.style={draw, rounded corners, minimum width=1.9cm,
        minimum height=0.9cm, align=center, font=\footnotesize},
      gd15emmi/.style={gd15zona, fill=red!15, draw=red!60!black},
      gd15save/.style={gd15zona, fill=blue!10, draw=blue!60!black},
      gd15freccia/.style={-{Stealth[length=2.2mm]}, thick},
      node distance=0.55cm]
    \node[gd15zona] (start) {Partenza\\(narrativa)};
    \node[gd15save, right=of start] (s1) {Salvataggio};
    \node[gd15emmi, right=of s1] (e1) {Zona EMMI\\(entra/esce)};
    \node[gd15save, right=of e1] (s2) {Salvataggio};
    \node[gd15emmi, right=of s2] (e2) {Zona EMMI\\(rientra/riesce)};
    \node[font=\footnotesize, right=0.35cm of e2] (dots) {$\cdots\ N$ volte};
    \draw[gd15freccia] (start) -- (s1);
    \draw[gd15freccia] (s1) -- (e1);
    \draw[gd15freccia] (e1) -- (s2);
    \draw[gd15freccia] (s2) -- (e2);
    \draw[gd15freccia] (e2) -- (dots);
    \draw[thick, gray] ($(e1.south west)+(0,-0.3)$) --
      node[below, font=\footnotesize, text=black, align=center, text width=8cm]
      {la zona EMMI taglia a metà ogni area: per proseguire bisogna
       attraversarla, e una volta visti si scappa verso l'uscita più vicina}
      ($(e2.south east)+(0,-0.3)$);
  \end{tikzpicture}
  \caption{Il ritmo di un'area di \emph{Metroid Dread} come lo descrive il
  docente: le zone EMMI sono la ``spina dorsale'' del livello, alternate ai
  salvataggi. Ricostruzione schematica, non la mappa del gioco.}
  \label{gd15:fig:emmi}
\end{figure}

\subsection{Riformulare: gli EMMI come pilastro e come USP}

Per tutte queste ragioni, dice il docente, \textbf{lo stealth non è un
pilastro, ma gli EMMI lo sono}: è più un problema di \textbf{riformulazione}
che un pilastro mancato, ``l'ordine è invertito''. Il fatto di avere un
\textbf{predatore inattaccabile} (fino a un certo punto, chiaramente), da cui
devi scappare, che pattuglia una zona specifica, e il fatto che esistano
queste zone, \emph{è} un pilastro del gioco. Anche la formulazione che
qualcuno ha usato, \textbf{``da predatore a preda''}, può essere un pilastro,
ed è esposta in maniera abbastanza convincente.

Non solo: questa minaccia degli EMMI, espansa, può diventare tranquillamente
un \textbf{unique selling point}. L'altro USP è ``senza ombra di dubbio''
un accorpamento di cose che molti studenti hanno scritto separatamente e che,
unite, fanno l'USP ideale del gioco:
\begin{itemize}
  \item è il \textbf{capitolo conclusivo} della saga di Metroid;
  \item è l'ennesimo capitolo \textbf{2D}, non legato alla serie
        \emph{Prime};
  \item è un \emph{metroidvania} di Nintendo su \textbf{Switch}.\footnote{Nella
        trascrizione ``Metroidvania di P-Nintendo su Switch''; ``P-'' è
        probabilmente un inciampo.}
\end{itemize}

\subsection{Lo stealth non è una meccanica}

Un inciso ``non rilevante, ma utile per il futuro'': lo
\textbf{stealth non è una meccanica}. Il \textbf{combattimento} è un
\textbf{sistema} composto da una serie di meccaniche, e rimane un sistema;
lo stealth è una \textbf{strategia}, un modo di affrontare la sfida (ci si
torna nella sezione~\ref{gd15:sec:cos-e}).

C'è anche un argomento storico a favore degli EMMI come pilastro. Dopo
\emph{Metroid Fusion} si è tentato più volte di fare un Metroid con al
centro \textbf{un predatore}, cioè un'intelligenza artificiale che dà la
caccia al giocatore (``se non erro, anche su DS''), e più volte i progetti
sono stati chiusi, a volte per limitazioni tecniche, a volte per i
costi.\footnote{Il docente lo riporta con cautela (``se non erro'',
``secondo me''). È noto che un \emph{Metroid Dread} era stato annunciato
internamente e mai realizzato all'epoca del Nintendo DS, e che l'idea del
predatore inarrestabile risale a \emph{Metroid Fusion} (con SA-X).} Poi è
uscito \emph{Dread}, che fa questa cosa in misura minore: non si parla di un
Nemesis\footnote{Il riferimento è al Nemesis di \emph{Resident Evil 3}, che
insegue il giocatore per la mappa.} che ti segue per la mappa e può entrare
più o meno ovunque. Gli EMMI sono confinati in zone ``per un discorso di
design e di semplificazione della produzione''. Anche questo li rende un
pilastro più della strategia stealth.

% ---------------------------------------------------------------------
\section{Il target audience e il doppio senso di ``estetiche''}
\label{gd15:sec:target}

\subsection{L'extra mile è benvenuto}

Il docente aveva chiesto di escludere alcune cose dalla vision (per esempio
la ``morale'', cioè il \emph{meaning}), ma è sempre contento se gli studenti
fanno l'\textbf{extra mile}: benvenga quindi chi ha aggiunto il
\textbf{target audience}. Va però detto che non era necessario, e che
l'unica cosa vista finora sul target è la \textbf{segmentazione numerica}
dell'utenza, per statistiche. Alberto conferma di aver diviso per età e
sesso, ``molto in linea generale'', senza entrare in categorie specifiche.
Ed è giusto così: nella costruzione di una vision quei campi fanno
riferimento a dati che vengono dalla \textbf{pre-produzione} (la ricerca e
la generazione di idee viste nei capitoli precedenti, capitolo~\ref{cap:ricerca-idee}),
e che qui mancano.

\subsection{Venerdì: i tipi di giocatori}

Quel campo della vision è l'applicazione di uno strumento che si inizia a
vedere \textbf{venerdì prossimo} (``per questo vi dicevo che è una lezione
importante''): a partire da quei dati si inizieranno a \textbf{separare tipi
di giocatori diversi} e ciò che chiedono in termini di emozioni e di
\emph{game aesthetics}. Francesca ha già fatto ricerca sugli
\textbf{archetipi} di base: si partirà da quelli e si espanderà.

È il primo passo per avere una \textbf{classificazione dell'utenza}, che
permette di prendere un'enorme mole di dati e approssimarla in grosse
categorie di giocatori. Nel tempo, fino al terzo anno (``è un percorso molto
lungo''), si arriverà a punti molto specifici di che cosa piace a un
giocatore e soprattutto di che cosa \textbf{non} gli piace. Questi giocatori
sono legati alla lista di emozioni già vista: ad alcuni piacerà soprattutto
l'\emph{excitement}, l'\emph{action} (come nelle motivazioni di Quantic
Foundry, sezione~\ref{sec:quantic-foundry}); altri saranno più legati alla
\emph{discovery}, che si ramifica a sua volta in \textbf{esplorazione} e in
\textbf{risoluzione di enigmi}. L'obiettivo è fare design \textbf{sempre
mirato a un'utenza}, sapendo perché le piacerà: se si inserisce un
combattimento action, sarà prediletto da un tipo di giocatore rispetto a un
altro, che magari preferisce un gameplay più strategico.

E non tutti i tipi di giocatori vanno accontentati allo stesso modo: per il
discorso della curva normale (sezione~\ref{gd10:sec:gaussiana}), ``il
vostro gioco non è per tutti'', quindi bisogna concentrarsi su
\textbf{uno}, al massimo \textbf{due} archetipi specifici.

\subsection{Attenzione alla parola ``estetiche''}

Nella vision gli studenti hanno usato ``estetiche'' nel senso delle
\emph{game aesthetics} legate alle emozioni; per ora va bene, ma per il
futuro il docente invita a fare attenzione. Nello stesso documento, o si usa
``estetiche'' \textbf{come Schell} (sezione~\ref{sec:schell}), cioè per
l'\textbf{aspetto} visivo e sensoriale, principalmente legato all'arte,
oppure la si usa per ciò che è legato alle \textbf{emozioni e alle
motivazioni} dell'utenza. Non mescolare i due sensi serve a evitare
confusione, ``non tanto per voi quanto per altri potenziali dipartimenti''.

% ---------------------------------------------------------------------
\section{Must do, must not, nice to have: il lato design}
\label{gd15:sec:must}

\subsection{Mea culpa: il mock-up fuorviante}

Il terzo punto è la parte più critica della revisione. Rifacendo le slide,
il docente ha mostrato nel \textbf{mock-up} di esempio della vision un blocco
di must do e must not che \textbf{non era il suo intento}: quelli mostrati in
primo piano erano legati alla \textbf{componente artistica} del gioco (le
texture, le mesh, il sound, tutto il \emph{dressing}), mentre quelli lato
design stavano più in basso. Averli messi davanti ha deviato l'attenzione,
e molti studenti hanno trattato più gli aspetti artistici che quelli
\textbf{meccanici e sistemici}. ``È colpa mia, mia colpa.'' ``Se fossi
artista vi darei un'altra risposta'', ma per un game designer la cosa che
interessa sono i must legati al \textbf{gameplay}.

\subsection{Che cosa ci interessa sapere}

Il lato design lo decide principalmente il dipartimento di design, o il
direttore creativo, a seconda di quanto è grande il progetto. E ciò che
interessa sapere è \textbf{quali meccaniche o quali sistemi ci sono per
forza}, con il maggior dettaglio possibile (pur con un certo margine di
creatività). Il docente prende un task a caso (quello di Francesco) e
costruisce gli esempi; la tabella~\ref{gd15:tab:must-dread} li raccoglie.

\begin{itemize}
  \item \textbf{Il melee counter.} Sappiamo che ci deve essere (Francesca
        l'ha messo nei pilastri), quindi va nei must do. Potremmo non
        sapere \emph{a quali} nemici si può fare counter e a quali no; qualcuno
        ha notato che i \textbf{boss} vengono tutti finiti con un QTE o un
        counter finale, cioè che il sistema si riproduce in forma ridotta
        anche lì, ma a quel punto della produzione potremmo non saperlo. Ci
        interessa sapere chi fa il counter, che eventualmente riprende dal
        capitolo precedente di MercurySteam;\footnote{Lo studio spagnolo che
        ha sviluppato \emph{Metroid: Samus Returns} (2017), in cui il melee
        counter è stato introdotto, e poi \emph{Metroid Dread}.} poi c'è
        margine per modificare le meccaniche o fare documenti a seconda di
        come sta uscendo il gioco.
  \item \textbf{La camera.} Prima della vision ci sarà stato un prototipo,
        quindi la camera è nota: molti hanno scritto \textbf{2.5D}. Camera
        laterale; i modelli 3D e gli sfondi in profondità sono più lato arte
        che design (``ve l'ho detto in maniera molto brutale e brutta''). Va
        scritto quale camera e quali movimenti.
  \item \textbf{Le abilità}, e se ritornano dai capitoli precedenti (magari
        non tutte). Il Phantom Cloak, pur non essendo un pilastro, è utile
        inserirlo qui; così come la Morph Ball, che c'è in tutti i Metroid
        e probabilmente ci sarà anche in questo.
  \item \textbf{Gli EMMI}: molti li hanno scritti nei pilastri e poi non li
        hanno ripresi nei must. Nei must do si scrive: un'intelligenza
        artificiale che insegue il giocatore, con una serie di
        \textbf{stati}, che non può essere uccisa se non a determinate
        condizioni. A quel punto della produzione potrebbe non essere chiaro
        che gli EMMI sono legati a delle zone, ma la vision si fa a
        posteriori, quindi si può segnare nei \textbf{must not}: gli EMMI
        \textbf{non possono uscire} dalle loro zone.
  \item \textbf{I boss}, chiaramente.
  \item \textbf{La mappa interconnessa}, divisa in una serie di
        \textbf{biomi}. Anche qui, a quel punto non si sa quanti. Alcuni
        hanno scritto il numero giusto: in quel caso si possono mettere
        ``tot biomi'' nei must do e un \textbf{``più o meno uno'' nei nice
        to have}. Se c'è più tempo di produzione si aggiunge un bioma, se ce
        n'è meno se ne toglie uno.
  \item \textbf{I sistemi accessori}, che quasi tutti hanno colto: senza di
        essi il gioco rimane lo stesso. L'\textbf{hard mode} (o \emph{dread
        mode}, ``come volete chiamarla''), i \textbf{collezionabili}, gli
        extra nel main menu: tutti tranquillamente togliibili.
\end{itemize}

\begin{table}[htbp]
  \centering
  \small
  \begin{tabular}{@{}p{0.31\textwidth}p{0.31\textwidth}p{0.31\textwidth}@{}}
    \toprule
    \textbf{Must do} & \textbf{Must not} & \textbf{Nice to have} \\
    \midrule
    Camera laterale (2.5D) e suoi movimenti &
    Gli EMMI non escono dalle loro zone &
    Un bioma in più o in meno ($\pm 1$) \\
    Movimento con locomozione specifica; combattimento melee e ranged &
    Nessuna valuta spendibile, nessun mercante, nessuno skill tree,
    nessun acquisto di oggetti &
    Hard mode / \emph{dread mode} \\
    Melee counter (eventualmente ereditato dal capitolo precedente) & &
    Collezionabili ed extra nel main menu \\
    Abilità, incluse quelle di ritorno (Morph Ball, Phantom Cloak) & &
    Puntine/segnalini colorati sulla mappa \\
    EMMI: IA cacciatrice a stati, uccidibile solo a certe condizioni & & \\
    Boss & & \\
    Mappa interconnessa a stanze, divisa in $N$ biomi; salvataggi e
    checkpoint, con un certo pacing & & \\
    Sistema di energia, missili & & \\
    \bottomrule
  \end{tabular}
  \caption{Esempi di must do, must not e nice to have \textbf{lato design}
  per \emph{Metroid Dread}, raccolti dagli esempi fatti a lezione. Non è una
  soluzione completa del task, ma la direzione indicata dal docente:
  sistemi e meccaniche, non texture e mesh.}
  \label{gd15:tab:must-dread}
\end{table}

\subsection{La mappa è un must, le puntine no}

L'esempio migliore della differenza tra must do e nice to have, però, sono
le \textbf{puntine}. Il docente cerca di mostrarle su \emph{Dread} (``input
database: no, non esiste''), poi ripiega su un altro gioco, ``tanto il
discorso è uguale'': chi ha giocato \emph{Hollow
Knight},\footnote{Nella trascrizione ``allo night'', ``Holonite'',
``Ullonite'', ``Allunite'', ``Olo Night''.} ma anche \emph{Dread}, ricorderà
che al giocatore vengono fornite \textbf{puntine, spilloni di colori
diversi} per segnare sulla mappa punti di interesse e tornarci, usabili con
una certa libertà, o \emph{agency}.
\begin{itemize}
  \item Avere la puntina in \emph{Dread} è un \textbf{nice to have}.
  \item Avere la \textbf{mappa}, fondamentale per navigare il mondo, è un
        \textbf{must do}. Un metroidvania senza mappa ``magari esiste,
        sicuramente esiste'', e può anche funzionare come leva di
        innovazione nel brainstorming con limitazioni; ma è problematico
        immaginare un Metroid in cui non puoi aprire la mappa, o accedervi,
        o sbloccarla in certi momenti.
\end{itemize}
Questa è la differenza, lato design di sistemi, tra i sistemi
\textbf{necessari}, da elencare nei must do, e quelli \textbf{accessori}, da
mettere nei nice to have.

\subsection{La valuta: un must not di Metroid}

In \emph{Hollow Knight} le puntine rientrano in un discorso più grosso, che
\emph{Metroid} non ha (e nemmeno una serie di Castlevania): la
\textbf{valuta spendibile di gioco}. In Metroid non c'è un mercante (come
alcuni hanno giustamente scritto, per il discorso dell'\textbf{isolamento}),
non puoi spendere nulla, non hai alberi delle abilità, non puoi comprare
oggetti: potrebbe tranquillamente essere un \textbf{must not}. In
\emph{Hollow Knight}, invece, è previsto.

Per arrivare al punto, il docente fa un sondaggio: chi ha giocato
\emph{Hollow Knight}? Solo Francesca (``non ditemi balle''). Allora:
\textbf{che cosa vende il mercante}? In \emph{Silksong}, ricorda una
studentessa, vende le puntine, vari tipi di segnalini (per i nemici, per i
punti di interesse, per differenziare mentalmente le cose), forse il
ciondolino per raggruppare le monete. Il docente aggiunge: anche in
\emph{Silksong} ci sono gli equivalenti dei \emph{charm} di \emph{Hollow
Knight}, passivi o attivi, e alcuni mercanti vendono anche le trappole.
E in \emph{Zelda}? Dopo un attimo di vuoto collettivo, Federica ricorda il
cibo da cucinare, qualcuno i semi per la fionda; il docente aggiunge le cose
``importantissime'': i \textbf{cuori} e i portacuori, le \textbf{borse}, il
potenziamento delle \textbf{bombe}, della \textbf{faretra} e così via. E in
\emph{Hollow Knight}/\emph{Silksong} alcuni mercanti vendono le
\textbf{chiavi}, i charm, il potenziamento della \textbf{maschera} (una vita
in più) o del \textbf{mana} (che in \emph{Silksong} è la seta).

\subsection{L'economy sink}

La domanda vera: davanti a un mercante che vende il portacuori e la puntina,
che cosa comprate? ``Se mi dite la puntina, siete dei matti.'' Una
studentessa obietta che dipende da quanti soldi si hanno: con pochi soldi, a
inizio gioco, forse prenderebbe la puntina per organizzarsi, sapendo che la
maschera resta lì. Interessante, risponde il docente, ma per ``buon uso e
costume di design'' il giocatore arriva al mercante avendo fatto una
moderata strada, con abbastanza soldi \textbf{non} per comprare una maschera
ma per arrivarci quasi. La scelta verte quindi sulla cosa più importante (la
maschera, la chiave); la puntina costa poco, ed è uno di quegli
\textbf{elementi accessori} che il giocatore compra quando ha \textbf{soldi
in eccesso}: perché ha già comprato tutto ciò che conta, o perché deve
spenderli da qualche parte (in \emph{Hollow Knight} e in \emph{Silksong} li
perdi morendo), o perché ne ha accumulati in quantità esagerata.

Non serviva quindi elencare proprio le puntine tra i nice to have (in quel
momento della produzione si potrebbe non sapere di volerle); ma in una vision
di \emph{Hollow Knight} ci sarebbe senz'altro, tra i nice to have, un qualche
tipo di \textbf{economy sink}: qualcosa nei mercanti che \textbf{assorbe il
denaro in eccesso} e lo ricambia con un oggetto di poco valore che però
\textbf{migliora l'esperienza}. L'economia di \emph{Hollow Knight} o di
\emph{Silksong} si regge in piedi anche senza gli oggetti accessori (``le
torrette che ti fanno le collane di palline''); certamente la migliorano.
(``Poi l'economia di \emph{Silksong} è un po' traballante, però\ldots
questo non lo devo dire io.'')

\paragraph{Il punto.} Il docente non voleva le puntine nel task: le ha usate
per mostrare la differenza. Per un designer, i \textbf{must do} della vision
sono soprattutto l'\textbf{elenco dei sistemi o delle meccaniche necessarie}
al gioco, e i \textbf{must not} tutto ciò che \textbf{non deve esserci}. E se
in un must do si scrive una cosa (per esempio ``mid poly''), non serve poi
ripetere nei must not il suo contrario (``niente high poly''): si spreca uno
spazio.

% ---------------------------------------------------------------------
\section{Fare ricerca su un gioco che non si conosce}
\label{gd15:sec:ricerca}

\subsection{L'osservazione di Federica}

Federica fa un'osservazione sul proprio lavoro: la cosa davvero difficile è
stata il tempo perso a \textbf{conoscere un gioco che non conosceva}; a
sfumature come i mercanti o le puntine non è neanche arrivata. ``Va
benissimo'', risponde il docente: le puntine erano un esempio, non una
richiesta. Ma ne approfitta per un discorso ``estremamente alto'', da
riprendere venerdì.

\subsection{Non capire non è un problema; non chiedere sì}

\textbf{Non è un problema non capire le cose; il problema è non chiedere},
non domandare, non migliorare. Il docente invita a confrontare il
\textbf{primo task consegnato} con quello di adesso: è presto per dirlo, ma
il miglioramento è notevole, e se si rifacesse ora una revisione sarebbe
completamente diversa, anche se è passata una settimana. Si va veloci per
necessità (tre lezioni a settimana, e gran parte del lavoro è a casa), e il
corso è iniziato da circa due mesi: ``c'è ancora tanto tempo, c'è ancora
tanta strada''. Non è una corsa ma una \textbf{maratona}, quella degli
studenti e quella del docente stesso, ``in generale nella disciplina che è il
game design, o il design in generale''. Quindi pazienza, \textbf{soprattutto
con se stessi}: chiedere ciò che non si è capito, ma anche darsi il tempo di
indagare, curiosare, farsi domande.

\subsection{La limitazione è voluta}

Non conoscere il gioco è \textbf{fatto apposta}, come per \emph{Pac-Man}: la
limitazione è voluta, la sfida è desiderata. In futuro capiterà spesso di
dover fare design su giochi o sistemi legati a titoli che non si conoscono;
quindi oltre a strutturare il documento c'è la \textbf{parte di ricerca}, e
una buona proporzione del tempo (``una buona metà'') va dedicata a quella, a
maggior ragione se il gioco è sconosciuto.

Ci saranno anche task a gioco libero, ``ci mancherebbe'', ma con un occhio
di riguardo: se il framework è adatto a un gioco basato sul level design e
gli studenti analizzano \emph{Doki Doki Literature Club}, o il gioco di
punizione dato a Valerio, ``quello dei piccioni''\footnote{Nella trascrizione
``Autoful Boyfriend'': è \emph{Hatoful Boyfriend}, un \emph{dating sim} in
cui i corteggiatori sono piccioni.} (un dating sim senza level design), lì
si pone il problema. Ma sono casi limite.

\subsection{Novanta per cento di ricerca}

Federica precisa: per lei il rapporto non è stato 50 e 50, ma \textbf{90 di
ricerca}, perché, pur avendo capito il concetto, non conoscendo le meccaniche
non sapeva cosa scrivere; guardava il gameplay di un paio di episodi e
annotava ciò che le sembrava interessante. Il docente risponde che conoscere
\textbf{più giochi possibile}, anche non solo gli ultimi usciti, dà più
spunti, non solo per fare giochi che piacciono o con cui si è confidenti,
ma per diventare designer migliori ``a 360 gradi''. E ciò che ha fatto
Federica è giusto.

\paragraph{Il metodo dei dieci secondi.} Per i must do e i must not, dato
una conoscenza di base del titolo, si può prendere il gioco a un certo
punto, guardarne un pezzo e \textbf{segnare le meccaniche} e ciò che succede.
Il docente lo fa in diretta su un video di \emph{Dread}: c'è
\textbf{movimento} con un sistema di locomozione specifico; c'è
\textbf{combattimento}, sia melee che ranged; c'è una \textbf{mappa} divisa
(un tempo a quadri, adesso a stanze), quindi un level design a
\textbf{stanze connesse} che compongono un bioma; un \textbf{sistema di
energia}, i \textbf{missili}, attacco pesante e attacco leggero. ``E abbiamo
visto veramente 10 secondi del gioco.''

Poi c'è una questione più profonda: capire a grandi linee \textbf{che tipo
di gioco} è. Molti, non conoscendolo, hanno faticato a dare il giusto peso
agli EMMI, perché \emph{Dread} è spinto in direzione horror lato marketing
(si chiama \emph{Dread}, dopotutto), ma \textbf{non è spiccatamente horror},
se lo si guarda nella sua interezza: anche questo può aver fuorviato.

\paragraph{Un consiglio pratico.} Per ottimizzare il tempo sui giochi che
non si conoscono, \textbf{dividere la qualità} dell'attenzione: momenti
dedicati, in cui la ricerca è profonda, e momenti in cui (``sarò brutale'')
si è in bagno, si sta mangiando, si sta facendo altro e si mette il gameplay
in sottofondo, annotando piccole cose sulle note del cellulare, su un foglio,
``sulla tovaglia'', che poi confluiscono nella ricerca principale. Con il
tempo si ottimizza anche questo.

% ---------------------------------------------------------------------
\section{Dividere i must per dipartimento, e la vestizione}
\label{gd15:sec:dipartimenti}

\subsection{La domanda di Francesca}

Francesca chiede: se la vision dev'essere una \textbf{visione globale} del
gioco, non è più intelligente fare must do e must not \textbf{divisi per
dipartimento}, level design, game design e così via? ``Infatti'', risponde
il docente: nel suo esempio la parte di design dei must è separata per
\textbf{dipartimenti di design}, con must do e must not legati al
\textbf{game design}, al \textbf{level design} e alla \textbf{narrativa};
tutto ciò che è legato all'\textbf{arte} sta altrove (ed è proprio il blocco
che è stato fuorviante). Non è l'unica soluzione, ma è quella che lui tende
a usare. La parte artistica di solito si desume dall'\textbf{art bible}
(che alcuni hanno citato) o dalle ricerche di direzione artistica, e confluisce
in must do e must not spiccatamente artistici. A lui, da designer, interessa
principalmente la parte di design.

\subsection{Sci-fi è vestizione}

Francesca precisa che nel suo task i must do non erano categorizzati
esplicitamente, ma si capiva che il primo era una visione globale, il
secondo di game design, il terzo di arte, il quarto di narrativa: è un
errore? ``Non è un problema'', ma il docente li rilegge.
\begin{itemize}
  \item \textbf{Sci-fi} è \textbf{vestizione}. Prendiamo un Castlevania,
        non i primi e nemmeno \emph{Symphony of the Night}, ma per esempio
        \emph{Aria of Sorrow},\footnote{Nella trascrizione ``Ariel of
        Sorrow''.} dove ci sono le armi ranged: togliamo il must do
        ``sci-fi'', e alcuni Metroid e alcuni Castlevania hanno
        \textbf{il 90\% del core in comune}. Certo, l'ambientazione si porta
        dietro aspetti di design (apre la possibilità di abilità
        futuristiche, la mimetizzazione tramite tecnologia, i tipi di raggio
        del cannone al braccio di Samus), ma resta soprattutto vestizione.
  \item \textbf{Dynamic combat} è chiaro: è un \textbf{sistema di
        gameplay}, e va bene.
  \item \textbf{Dark environment} è un discorso spiccatamente
        \textbf{artistico}: ambienti e zone tetre, il \emph{mood}. Può
        legarsi al level design (quanto grandi o piccoli sono i corridoi e
        le stanze), ma ha più conseguenze sull'arte.
  \item \textbf{Narrativa}: comprensibile.
  \item \textbf{Texture} e \textbf{sound} sono lato arte: da una parte il
        3D, dall'altra il sound design (che ha ``design'' nel nome, ma è una
        disciplina più artistica).
\end{itemize}
Dunque solo il secondo punto è legato al game design. Gli altri non sono
sbagliati: \textbf{restringono il campo} su che cosa dovrebbe essere
\emph{Metroid Dread}. Ma servono \textbf{più elementi} per definire che cosa
rende \emph{Dread} quel gioco lì \textbf{con il controller in mano}: una
volta che il giocatore muove il personaggio, come lo vede, che camera ha.

\subsection{Tre Zelda, tre camere}

Per mostrare quanto conta la camera il docente mette a confronto tre
\emph{Zelda} (tabella~\ref{gd15:tab:zelda}): \emph{A Link to the Past},
\emph{A Link Between Worlds} e un Zelda moderno, \emph{Breath of the Wild}.
I primi due hanno camere simili, eppure appartengono a epoche completamente
diverse; il secondo e il terzo appartengono alla stessa epoca, ma con camere
completamente diverse. \emph{Breath of the Wild} ha la classica terza
persona, con un braccio della camera molto più lungo, cioè una camera che
sta molto più lontana dal personaggio: perché è un \textbf{open world} e un
gioco \textbf{emergente} (se ne parlerà). Il design lascia al giocatore
un'enorme quantità di possibilità (manipolare mostri, far esplodere eventi a
catena, incendiare alberi), e quindi quella camera ha senso.

\begin{table}[htbp]
  \centering
  \small
  \begin{tabular}{@{}p{0.27\textwidth}p{0.3\textwidth}p{0.35\textwidth}@{}}
    \toprule
    \textbf{Gioco} & \textbf{Camera} & \textbf{Osservazione del docente} \\
    \midrule
    \emph{A Link to the Past} & Top-down / isometrica ``classica'' &
      Simile alla successiva, ma di un'epoca completamente diversa. \\
    \emph{A Link Between Worlds} & Top-down, con momenti a camera laterale
      (si diventa bidimensionali, ``come in \emph{Mario Odyssey}'') &
      Stessa epoca di \emph{Breath of the Wild}, camera completamente
      diversa. \\
    \emph{Breath of the Wild} & Terza persona con braccio lungo, molto
      lontana dal personaggio & Open world ed emergente: la camera ampia
      serve a un design che lascia fare moltissime cose. \\
    \bottomrule
  \end{tabular}
  \caption{Le tre camere di Zelda usate per mostrare perché la camera va nei
  must do. Le classificazioni sono quelle date a voce dal docente, che
  stesso esita (``isometrica, no? Beh, oddio'').}
  \label{gd15:tab:zelda}
\end{table}

Poi il \textbf{personaggio}: chi è, che cosa può permettersi di fare, quali
\textbf{abilità} ha. Poi il \textbf{level design}: gli ambienti sono connessi
o no? Esiste il \emph{fast travel} o no? Sono questi elementi a comporre
l'anima del gioco a livello di design.

\subsection{Il test: immaginate il gioco dalla vostra lista}

Il test che il docente propone: \textbf{leggete la vostra lista} di must do
\textbf{senza pensare al gioco}, e provate a creare nella vostra testa un
gioco con i punti che avete descritto. \textbf{Quanto si avvicina a
Metroid?} Se si allontana, o se \emph{può} allontanarsi notevolmente,
allora c'è un problema. Nella vision bisogna avere una linea
\textbf{estremamente netta} di che tipo di genere (``discorso strano,
eh''), che tipo di estetiche e che tipo di meccaniche ha l'esperienza; più
si va nello specifico, meglio è.

\subsection{Il numero di abilità e i costi}

Con \emph{Dread}, un gioco già uscito, gli studenti potevano analizzare
tutto; in una vision vera la precisione assoluta, per esempio sul numero di
abilità, non c'è. L'obiettivo però è arrivarci: in \emph{Breath of the Wild}
hai il \textbf{magnete}, le \textbf{bombe}, l'\textbf{aliante}, l'\textbf{arco}
e così via. Se si riesce a definire il \textbf{numero minimo di oggetti e di
abilità} che serve, si ha già chiaro \textbf{quanto costerà la produzione}:
non ci si può permettere 20 abilità con un anno di sviluppo o con un budget
estremamente ridotto. Quindi: immaginare il gioco con gli elementi minimi,
inserire nella vision la lista di ciò che rende quel gioco tale, e poi tutte
le cose accessorie che, ``se ci sono, ben venga; se non ci sono, pace''.

\subsection{Lo skill tree: must do o nice to have?}

Lo stesso sistema può stare in categorie diverse a seconda del gioco. Lo
\textbf{skill tree} è un nice to have in alcuni giochi e un must do in altri;
il docente chiede esempi, e arrivano.
\begin{itemize}
  \item In \emph{Borderlands} lo skill tree è un \textbf{must do}.
  \item In un gioco di \emph{Star Wars} mostrato a schermo è meno
        rilevante.\footnote{Il titolo non è nominato; con ogni probabilità
        uno dei \emph{Star Wars Jedi}, ma non è verificabile dalla fonte.}
  \item In \emph{League of Legends} le \textbf{rune} sono lo skill tree del
        gioco, un sistema cambiato tante volte. Serve a dare
        \textbf{agency} al giocatore per modificare \textbf{a monte} lo stile
        di gioco a parità di personaggio. L'esempio: il docente e Daniele
        giocano lo stesso campione, Garen; il docente lo vuole più
        \textbf{tank}, Daniele più \textbf{DPS}. Serve un modo per
        personalizzare il personaggio prima di fare la build in partita (una
        personalizzazione non estrema).
  \item E \emph{Batman}?\footnote{Probabilmente uno degli \emph{Arkham};
        la fonte non lo precisa.} Anche lì bisognerebbe andare a vedere.
\end{itemize}
La domanda da porsi è sempre: \textbf{il gioco sarebbe lo stesso con o
senza quel sistema? Quanto si allontana?} ``Questa è la domanda da un milione
di dollari a cui dovete rispondere nella vision.''

% ---------------------------------------------------------------------
\section{Gli esempi non sono la regola}
\label{gd15:sec:esempi}

Mentre gli studenti pensano ad altre domande, il docente porta una questione
sul \textbf{compendium}: l'esempio che aveva mostrato non va seguito uno a
uno (lo ha anche scritto nelle slide). Era un pezzo del compendium di una
\textbf{produzione estremamente grossa}; uno studente potrebbe non avere un
elemento ``generale'', potrebbe chiamarlo ``visione'' o in un altro modo,
potrebbe spostare elementi da un blocco all'altro in base a come vuole
organizzare la produzione. Nel progetto annuale, per esempio, il blocco
\textbf{marketing} probabilmente non interesserà, almeno non subito, e
finché non è utile si può ignorare. Se invece lo si inserisce per
sperimentare, per curiosità su che cosa si trova in un dipartimento di level
design o di che cosa si preoccupano gli artisti o i sound designer, è un
altro discorso.

L'invito è a \textbf{non restare ancorati agli esempi} dove c'è margine di
manovra, e a \textbf{sperimentare} il più possibile (entro limiti ragionevoli).
Gli studenti, ``più di me'', hanno una mente ancora aperta; col tempo c'è
il rischio, ``e va bene'', di chiudersi sulle cose che funzionano. Più si
prova adesso, più si impara, e magari si \textbf{innova}, usando un framework
in modo nuovo: il docente racconta di aver visto, tra gli studenti,
``cose che voi umani non potete immaginare'', cioè usi di framework a cui lui
stesso non aveva pensato e che si sono rivelati estremamente utili. Nel
dubbio (seguire nettamente o è una linea guida?), chiedere: ``alla peggio, se
non vi rispondo, è perché c'è ancora più margine''.

% ---------------------------------------------------------------------
\section{Il task dei feedback: stesso mandate, risposte diverse}
\label{gd15:sec:zoe}

L'ultimo punto della revisione riguarda il task dei \textbf{feedback} sulle
concept art dello zombie game (sezione~\ref{gd10:sec:task-zoe}). Era un task
in cui designer diversi partivano dallo \textbf{stesso mandate}: lo zombie
game, la segmentazione dell'utenza, le estetiche di gioco e le emozioni
scaturite, la camera, il budget erano uguali per tutti. Eppure molti hanno
dato feedback \textbf{molto lontani} sullo stesso concept (alcuni comuni, e
ci sta). È la parte, ``quel 50\% di creatività del lavoro'', che il docente
vuole che gli studenti guardino subito: \textbf{come gli altri hanno risposto
in modo diverso allo stesso input}. Nel dare e nel ricevere feedback si
possono ricevere cose molto diverse, e sta al designer domandarsi qual è la
soluzione migliore a seconda del mandate, della visione e di ciò che vuole
fare.

L'esempio: magari un nemico non lo si vuole, perché il \textbf{braccio a
frusta} è estremamente costoso. Chi lo pensa è un tipo più coscienzioso, che
preferisce una soluzione \textbf{sicura} a una \textbf{innovativa e
rischiosa}; non è un problema. Altri punteranno all'innovazione, nonostante
i rischi. Nessuna delle due risposte è sbagliata: lo si vede solo in
\textbf{produzione}, e qui i soldi veri non c'erano e il gioco nemmeno. Anche
il docente non dice ``avete sbagliato'': è d'accordo con alcuni punti e in
disaccordo con altri, ma chi ha ragione lo si vedrebbe solo facendo
``questo incredibile zombie game, cosa che ahimè non possiamo fare''.
L'invito è ad andare a leggere i commenti degli altri.

% ---------------------------------------------------------------------
\section{Un flowchart da fare insieme (rimandato)}
\label{gd15:sec:flowchart}

\subsection{Il gioco: Grotte}

Nessun'altra domanda emerge (``non vi è bollito nulla dal cervello''), e il
docente annuncia il programma per l'ultima mezz'ora: fare insieme il
\textbf{flowchart che descrive l'intero sistema di gioco} di
\textbf{Grotte}.\footnote{Il nome è riportato così nella trascrizione; non
è stato possibile verificare con certezza di quale gioco si tratti. Dalla
descrizione è un piccolo gioco di ruolo OSR, con regole riassunte nelle slide
del docente e un manuale completo a pagamento (``7 euro''). Il docente
chiede anche chi conosca ``Modborg'', quasi certamente \emph{Mörk Borg},
un altro gioco di ruolo di stampo OSR; ma il gioco su cui fare il flowchart
è ``Grotte''.} Chiede chi conosca \emph{Mörk Borg} e chi giochi a
\textbf{giochi di ruolo} in generale (una domanda già fatta alla seconda
lezione da Jack e Diego): è un buon modo per stringere il legame tra
compagni, dedicando delle sessioni (``mi piacerebbe partecipare io stesso,
ma ahimè ci separa il ruolo docente e studente''). Grotte è un punto di
partenza interessante perché è un \textbf{OSR}\footnote{\emph{Old School
Renaissance} (o \emph{Revival}): giochi di ruolo che recuperano lo spirito
delle prime edizioni di \emph{Dungeons \& Dragons}, con regole leggere e
molta improvvisazione del master.} con \textbf{regole estremamente
semplificate}: si può mappare tutto, e si gioca con pochissimo, carta e
qualche dado a sei facce. Le regole sono nelle slide; quelle complete si
trovano online, oppure si possono comprare (``7 euro da dare a questo
cristiano''). Ma l'obiettivo è \textbf{mappare l'intero sistema} con un
flowchart.

\subsection{I due errori più comuni}

Il docente anticipa i due problemi principali visti nei task (riprendendo le
regole della sezione~\ref{sec:flowchart-regole} e i blocchi della
sezione~\ref{sec:flowchart-blocchi}), schematizzati nella
figura~\ref{gd15:fig:errori-flowchart}:
\begin{enumerate}
  \item \textbf{blocchi slegati}: un blocco che non va da nessuna parte o
        che non ritorna nel flusso;
  \item \textbf{blocchi azione con uscite multiple}: usare il rettangolo
        con più frecce in uscita è \textbf{sbagliato}; nelle regole base del
        flowchart, se si vuole una ramificazione, c'è il blocco apposito, il
        \textbf{rombo} della decisione.
\end{enumerate}

\begin{figure}[htbp]
  \centering
  \begin{tikzpicture}[
      gd15fc/.style={draw, thick, minimum height=0.75cm, align=center,
        font=\footnotesize},
      gd15az/.style={gd15fc, rectangle, minimum width=2.1cm},
      gd15dec/.style={gd15fc, diamond, aspect=1.9, inner sep=1pt},
      gd15ter/.style={gd15fc, rounded rectangle, minimum width=1.8cm},
      gd15f/.style={-{Stealth[length=2mm]}, thick},
      node distance=0.55cm and 0.6cm]
    % Errore 1
    \node[font=\footnotesize\bfseries] (t1) at (0,0) {Errore 1: blocco slegato};
    \node[gd15ter, below=0.35cm of t1] (a1) {Inizio};
    \node[gd15az, below=of a1] (b1) {Tira il dado};
    \node[gd15az, below=of b1] (c1) {Fine turno};
    \node[gd15az, draw=red!70!black, left=0.4cm of b1] (x1) {Cura};
    \draw[gd15f] (a1) -- (b1);
    \draw[gd15f] (b1) -- (c1);
    \node[font=\scriptsize, text=red!70!black, below=0.1cm of x1,
      align=center] {né entra\\né esce};
    % Errore 2
    \node[font=\footnotesize\bfseries] (t2) at (4.3,0) {Errore 2: due uscite};
    \node[gd15az, below=0.35cm of t2] (b2) {Attacca};
    \node[gd15az, below left=0.7cm and -0.4cm of b2] (c2) {Colpito};
    \node[gd15az, below right=0.7cm and -0.4cm of b2] (d2) {Mancato};
    \draw[gd15f, red!70!black] (b2) -- (c2);
    \draw[gd15f, red!70!black] (b2) -- (d2);
    % Corretto
    \node[font=\footnotesize\bfseries] (t3) at (8.9,0) {Corretto: il rombo};
    \node[gd15az, below=0.35cm of t3] (b3) {Attacca};
    \node[gd15dec, below=of b3] (q3) {Colpito?};
    \node[gd15az, below=0.6cm of q3] (s3) {Danno};
    \node[gd15az, right=0.45cm of q3] (n3) {Nulla};
    \draw[gd15f] (b3) -- (q3);
    \draw[gd15f] (q3) -- node[left, font=\scriptsize] {sì} (s3);
    \draw[gd15f] (q3) -- node[above, font=\scriptsize] {no} (n3);
  \end{tikzpicture}
  \caption{I due errori sui flowchart segnalati dal docente (a sinistra e al
  centro) e la forma corretta della ramificazione (a destra), con uscite del
  rombo etichettate e il ``sì'' in basso. Esempio inventato a scopo
  illustrativo.}
  \label{gd15:fig:errori-flowchart}
\end{figure}

Il flowchart vero e proprio, però, il docente lo vuole fare dopo aver
``recuperato'' la lezione, cioè dopo la parte nuova: e, come si vedrà, il
tempo non basterà (sezione~\ref{gd15:sec:prossima}).

% ---------------------------------------------------------------------
\section{Il calendario e il task di pitching}
\label{gd15:sec:calendario}

\subsection{Le date}

Una comunicazione ``importantissima'', da girare anche agli assenti (le
domande e le discussioni restano per venerdì):\footnote{Le date sono
riportate come le dice il docente nella trascrizione (``venerdì prossimo,
27'', poi ``il 14'', ``il 18'' e ``il 21''), senza mese esplicito. Il 27
coincide con il venerdì 27 novembre citato anche per il Beer to Peer; il 14,
il 18 e il 21 sono con ogni probabilità date di dicembre, ma la sequenza non
è del tutto coerente nella trascrizione e va verificata sul calendario del
corso.}
\begin{itemize}
  \item c'è lezione \textbf{venerdì prossimo, il 27};
  \item poi \textbf{si salta il 14}: quel giorno gli studenti avranno
        lezione con \textbf{Diego e Jack}, da remoto, e il docente non potrà
        partecipare;
  \item la lezione successiva del docente è quella del \textbf{18}.
\end{itemize}

\subsection{Un task di gruppo da presentare in aula}

La lezione del 14 sarà dedicata a \textbf{come presentare}, cioè al
\textbf{pitching}: come farlo, i tipi di pitch e così via. In quella lezione
verrà dato un task (``che sorpresa''), e normalmente il docente direbbe di
consegnarlo per il 21 e lo commenterebbe quando gli studenti vogliono; ma il
21 non ci si vede, ci si vede il 18. Ha provato ``in mille modi'' a spostare
le lezioni (lui potrebbe il 19 o il 23, ``ma probabilmente voi no''), e
questo è il miglior caso possibile. Quindi:
\begin{itemize}
  \item il task va \textbf{consegnato per forza il 18};
  \item sarà un \textbf{task a gruppi}: il 14 verrà data una
        \textbf{presentazione da fare}, e il 18 la si farà \textbf{in
        classe};
  \item i \textbf{feedback} saranno \textbf{a voce}, sulla
        \textbf{performance in aula}, non su un documento commentato.
\end{itemize}
Il docente chiede di tenersi liberi il 18, per venire in presenza e per
lavorare in gruppo nei pochi giorni a disposizione (quattro, cinque contando
il lunedì): è importante imparare a farlo nel più breve tempo possibile,
prima della \textbf{pausa invernale}.\footnote{Nella trascrizione ``pausa
internale''.} Lo dice con largo anticipo proprio perché il tempo sarà poco.

Prima della pausa il docente chiude la condivisione (per non far vedere agli
studenti l'anteprima della slide successiva, ``ed è un casino'') e concede
sei minuti: ``ci vediamo alle 40'', ``freschi come delle belle rose''.

\subsection{Il Beer to Peer}

Al rientro, in attesa degli altri, un'altra comunicazione (forse Daniele
l'aveva mandata solo sul Discord di narrativa): il \textbf{27 novembre} c'è
il \textbf{Beer to Peer}\footnote{Nella trascrizione ``peer-to-peer'', poi
``Beer to Peer di GDA'' e ``Virtual Pier GDA''. È un incontro informale tra
sviluppatori a Milano; l'organizzazione indicata come ``GDA'' non è
identificabile con certezza dalla fonte.} a Milano: un momento di incontro
con sviluppatori, in un posto che cambia ogni volta, alle 19. Il docente non
sa se riuscirà ad andare (``magari dopo cena''). È presto, è il primo anno,
ma è già un buon modo per conoscere persone, chiacchierare con sviluppatori
e fare \textbf{networking}.

% ---------------------------------------------------------------------
\section{Dopo la pausa: le tre porte}
\label{gd15:sec:porte}

\subsection{Scegliete una porta}

La lezione nuova parte con ``l'incredibile domanda'': sullo schermo ci sono
\textbf{tre porte}, una \textbf{gialla}, una \textbf{rossa} e una
\textbf{blu} (figura~\ref{gd15:fig:porte}). Quale scegliete? Una studentessa
sceglie la gialla (``preferisco il giallo rispetto al rosso e al blu''),
nessuno la rossa, tutti gli altri la blu, ``per esclusione''. Federica,
rientrata, conferma la blu (ha appena iniziato \emph{Blue Prince}:
``da quel poco che ho visto, gran gioco'').\footnote{Nella trascrizione
``ho appena giocato a Blueprints''; dal contesto (porte, scelta, ``gran
gioco'') è quasi certamente \emph{Blue Prince}, un gioco di esplorazione in
cui si scelgono stanze dietro le porte di una villa.}

Il docente apre le porte. ``Non è il trucco'' (quello del presentatore che
apre una porta vuota e chiede se volete cambiare scelta per aumentare le
probabilità):\footnote{Il riferimento implicito è al problema di Monty Hall.}
il grande sacco di soldi era dietro la \textbf{porta rossa}. Tutti gli altri
si beccano le \textbf{tigri del Bengala}, affamate, ``che devoreranno le
vostre carni''. È, scherza, l'unico caso in cui tutti si sono presi le tigri.

\begin{figure}[htbp]
  \centering
  \begin{tikzpicture}[
      gd15porta/.style={draw, thick, minimum width=1.6cm, minimum height=2.6cm,
        rounded corners=2pt}]
    \node[gd15porta, fill=yellow!45] (g) at (0,0) {};
    \node[gd15porta, fill=red!45] (r) at (2.6,0) {};
    \node[gd15porta, fill=blue!35] (b) at (5.2,0) {};
    \foreach \n in {g,r,b} {\fill (\n.east) ++(-0.3,0) circle (0.07);}
    \node[below=0.15cm of g, font=\footnotesize, align=center] {gialla\\
      \textbf{tigri}};
    \node[below=0.15cm of r, font=\footnotesize, align=center] {rossa\\
      \textbf{sacco di soldi}};
    \node[below=0.15cm of b, font=\footnotesize, align=center] {blu\\
      \textbf{tigri}};
    \node[above=0.2cm of r, font=\footnotesize, align=center, text width=8cm]
      {Prima di scegliere: nessuna informazione sulle porte,\\
       nessuna informazione sulla ricompensa (o sulla punizione)};
  \end{tikzpicture}
  \caption{La scelta delle tre porte con cui si apre la lezione: tre opzioni,
  come vuole la regola dell'agency (capitolo~\ref{gd9:cap:agency-visione}),
  ma senza alcuna informazione né prima né sulla posta in gioco.}
  \label{gd15:fig:porte}
\end{figure}

\subsection{Che cosa non va in questa scelta}

La scelta rispetta i limiti visti nella lezione sull'agency
(sezione~\ref{gd9:sec:troppe-scelte}): non un milione di porte, ma
\textbf{tre}, esattamente come cioccolato, gelato e torta di mele (``buona
la torta di mele''). ``Tuttavia c'è un ma.'' Siete soddisfatti della vostra
scelta? Nessuno risponde (``beh, comunque è una buona morte, dai''): siete lì
con la tigre che vi divora la giugulare mentre il presentatore apre la porta
con un milione di dollari, ``gioco distopico al 100\%''. Al netto di
``quanto exciting sia morire per zampa di una tigre del Bengala'', perché
questa scelta, implementata in un videogioco, \textbf{non sarebbe l'ideale}?

Una studentessa propone una lettura di simbologia: il \textbf{rosso}, anche
nei cartelli stradali, è comunemente associato al \textbf{pericolo}; nello
\emph{Star Wars} il blu ne è la controparte; il giallo magari è un bonus
(``ma sto sparando a caso''). È una riflessione interessante legata
all'\textbf{affordance}, che si affronterà in futuro, soprattutto con
Stefano. Ma il docente la riporta al punto: la domanda è fatta \emph{a
posteriori}, dopo aver aperto le porte. \textbf{Prima} di aprirle gli
studenti non avevano informazioni:
\begin{enumerate}
  \item \textbf{nessun indicatore} su quale porta fosse quale, come invece
        capita nei giochi (anche in \emph{The Stanley Parable}): quale porta
        è accessibile e quale no, dove concentrare l'interesse, segnali di
        pericolo chiaramente leggibili o fraintendibili. Nessun NPC aveva
        detto ``state attenti alla porta rossa'', né c'era il grande tranello
        del ``il rosso è il colore del pericolo, ma in questo caso no'';
  \item \textbf{nessuna informazione sulla ricompensa}: né sulla sua
        grandezza né sul suo contenuto (``qualcuno di voi magari è
        milionario''), né su che cosa sarebbe arrivato dopo, ricompensa o
        punizione.
\end{enumerate}
Una scelta limitata, come vuole l'agency, ma fatta \textbf{senza
informazioni prima} e \textbf{senza sapere a che cosa si andava incontro
dopo}. È il secondo passo del discorso: dopo la \textbf{libertà} di agency,
\textbf{come si progetta una scelta che venga percepita come
\emph{meaningful}}.

% ---------------------------------------------------------------------
\section{Che cos'è una decisione}
\label{gd15:sec:cos-e}

\subsection{Dalla colazione al finale}

Per capire che cosa fa una decisione dentro un gioco, il docente mostra
alcuni esempi.
\begin{description}
  \item[\emph{Life is Strange}, la colazione] Una scelta estremamente
        minore, in una cutscene:\footnote{Nella trascrizione ``Cazzin'' e
        ``Kazin'' per \emph{cutscene}, qui e oltre.} ti chiedono che cosa
        vuoi per colazione, \textbf{pancake} o \textbf{uova e bacon}.
        Qualsiasi scelta tu faccia cambia poco ai fini della
        trama.\footnote{Dal racconto successivo (la protagonista che si
        sveglia, va dall'amica, vanno a mangiare con la madre) la scena
        sembra appartenere a \emph{Life is Strange: Before the Storm}
        piuttosto che al primo capitolo; il docente la attribuisce a
        \emph{Life is Strange} in generale.}
  \item[\emph{BioShock}, le Sorelline] In netta contrapposizione. Le
        \textbf{Little Sister} sono bambine orfane in cui è stata impiantata
        una lumaca parassita che le trasforma in una sorta di abominio,
        pur mantenendo la loro umanità; sono un mezzo estremamente efficiente
        per produrre l'\textbf{ADAM}, il ``mana'' di \emph{BioShock}. Il
        giocatore sceglie se \textbf{terminare le loro sofferenze} e
        guadagnare di conseguenza, oppure \textbf{salvarle} togliendo la
        lumaca. La scelta non solo è percepita come nettamente diversa da
        ``uova e bacon'', ma porta a \textbf{finali diversi}: secondo il
        docente, best ending se le hai salvate tutte, bad ending (con un
        leggero twist) se ne hai salvate solo alcune, bad ending se le hai
        uccise tutte, prosciugandole dell'ADAM.\footnote{Così nella
        trascrizione. Nel gioco il finale ``buono'' richiede di non
        prosciugarne nessuna, e i finali negativi hanno varianti a seconda
        di quante se ne sono prosciugate.}
  \item[\emph{Fallout 3}, Megaton] La scelta di cui si parla spesso (ne
        aveva parlato anche Diego nella lezione condivisa): \textbf{non
        bombardare} con un'atomica una città i cui abitanti ti hanno trattato
        male, oppure \textbf{farla saltare} e ottenere un sacco di soldi. Il
        ragionamento di Diego: ``una nucleare in più, una nucleare in meno,
        nel mondo di \emph{Fallout} che male fa?''. Ma, considerando la
        componente etica, stai sterminando una città.
  \item[\emph{Life is Strange}, le grandi scelte] Quelle che ramificano un
        finale possibile.
\end{description}
Tutte rientrano nello spettro delle \textbf{decisioni}.

\subsection{Tutto è una decisione}

``Vi dirò di più'': a conti fatti, \textbf{qualsiasi cosa il giocatore fa}
in un gioco è una decisione. Andare a destra o a sinistra, cioè premere
l'input a destra o a sinistra, è una decisione. Raggiungere una città in
300 passi invece che in 301, in un gioco a tile come \emph{Pokémon}, è una
decisione (che per semplicità si ignora).

Lo è anche \textbf{con che cosa si gioca}: la tecnologia usata, cioè il
controller, cambia l'approccio al gioco (più o meno efficiente nella mira con
mouse e tastiera, più o meno efficiente nei combattimenti dei Souls,
negli action o nei platform con il controller). Nei \textbf{picchiaduro} (``è
Francesco il maestro di \emph{Smash Bros.}, ma oggi non c'è'') qual è la
periferica ideale per la competizione ad alto livello? Il docente mostra un
video di un torneo (``Evo, Tekken, 2025''): si usano gli \textbf{arcade
stick}, con leva e bottoni, come nei cabinati. Per chi non è avvezzo può
essere sorprendente, ma gran parte dei picchiaduro a livello professionistico
si gioca così (e non costano poco): nell'apprendimento delle combo e
nell'efficienza in gioco dà un vantaggio notevole, l'equivalente di mouse e
tastiera negli sparatutto. (C'è anche il \textbf{Power Glove} di Nintendo,
``per farsi una risata''.) Ci sono framework che considerano proprio i
\textbf{limiti} che il giocatore si impone o che gli si danno; se ne
riparlerà. Questa è una decisione che avviene \textbf{fuori} dal gioco, una
\textbf{meta-decisione}, ma è parte dei tipi di decisione possibili.

Sono decisioni, logicamente, anche la \textbf{classe} che si sceglie, lo
\textbf{spread delle statistiche}, dove investire i punti, e \textbf{come
affrontare i boss}.

\subsection{Lo stealth come strategia}

E qui il docente si riallaccia allo stealth
(sezione~\ref{gd15:sec:pilastri}): lo stealth è una \textbf{strategia},
perché è legato a \emph{come} affronto la sfida, cercando di non essere
visto. In molte missioni di \emph{Metal Gear} c'è l'esito \emph{ghost} se
non vieni visto da nessuno, e puoi giocarle completamente in stealth; ma
puoi anche sterminare tutti. Lo si vede ancora meglio in \emph{Hitman}: la
\textbf{prima missione} di \emph{Hitman World of Assassination}, ``se non
erro'', è gratuita, e il docente consiglia di giocarla proprio per le
decisioni. Le possibilità con cui eliminare il bersaglio, il prendere un
oggetto e lanciarlo contro un altro, l'osservare le ronde delle persone o
che cosa fanno dentro e fuori dal turno di lavoro: tutto dà un'enorme
\textbf{agency}. Non a caso è un \textbf{sandbox}, ``un grosso parco giochi''
in cui raggiungere l'obiettivo in una quantità di modi esagerata; ed è uno
degli unique selling point di \emph{Hitman}.

\paragraph{Che cosa ci interessa.} La prima cosa da considerare, dunque, è
che \textbf{qualsiasi cosa il giocatore fa} in un gioco, \textbf{e il modo
in cui vi si approccia}, è una decisione (a volte anche fuori dal gioco).
In analisi e in design, però, interessano soprattutto le decisioni che si
possono \textbf{controllare o manipolare}: come raggiunge una porta, quale
porta sceglie, che classe usa rispetto al tipo di giocatore che è, come
affronta un boss.

\paragraph{Approfondimento: un gioco è una serie di decisioni.} È la stessa
premessa di \emph{Challenges for Game Designers} di Brenda Brathwaite e Ian
Schreiber: ``distillato alla sua essenza, il game design consiste nel creare
per i giocatori opportunità di prendere \textbf{decisioni significative} che
influenzano l'esito del gioco''. Gli esempi del libro spaziano dal
posizionamento delle truppe negli strategici all'allocazione dei punti nei
GDR, dal pezzo da muovere negli scacchi al mirare e sparare in un FPS, fino a
premere il tasto giusto al momento giusto in \emph{Guitar Hero}: anche
l'azione puramente fisica è una decisione. E quando non c'è scelta (il
\emph{Monopoli} a proprietà tutte vendute, ``tira i dadi e paga''), il
gioco si spegne.

% ---------------------------------------------------------------------
\section{I tipi di decisione}
\label{gd15:sec:tipi}

Per ora interessa \textbf{categorizzare} le decisioni. Il docente mette ``le
mani avanti'': le categorie \textbf{non sono mutualmente esclusive}. Una
decisione può appartenere a più tipi contemporaneamente: essere legata al
tema e alle meccaniche allo stesso modo, oppure avere una risoluzione
soprattutto meccanica con elementi di tema. Le categorie sono tre, con le
loro sottocategorie (figura~\ref{gd15:fig:tipi}).

\begin{figure}[htbp]
  \centering
  \begin{tikzpicture}[
      gd15rad/.style={draw, thick, rounded corners, fill=gray!15,
        font=\small\bfseries, minimum width=2.6cm, minimum height=0.8cm},
      gd15cat/.style={draw, thick, rounded corners, fill=blue!8,
        font=\footnotesize\bfseries, minimum width=2.8cm, minimum height=0.7cm},
      gd15sub/.style={draw, rounded corners, font=\scriptsize, align=left,
        text width=3.3cm, inner sep=3pt},
      gd15l/.style={thick, -{Stealth[length=1.8mm]}}]
    \node[gd15rad] (d) at (0,0) {Decisioni};
    \node[gd15cat] (t) at (-4.6,-1.4) {Tema};
    \node[gd15cat] (m) at (0,-1.4) {Meccaniche};
    \node[gd15cat] (x) at (4.6,-1.4) {Meta};
    \foreach \c in {t,m,x} \draw[gd15l] (d) -- (\c);
    \node[gd15sub, below=0.35cm of t] (t1) {\textbf{Emotion driven}:
      viscerali, d'istinto (Megaton)};
    \node[gd15sub, below=0.2cm of t1] (t2) {\textbf{Narrative e sociali}:
      finali, linee narrative, legami, romance (\emph{Hades})};
    \node[gd15sub, below=0.2cm of t2] (t3) {\textbf{Estetiche}: dove esploro,
      \emph{Fashion Souls}, le build di \emph{Monster Hunter}};
    \node[gd15sub, below=0.35cm of m] (m1) {\textbf{Strategia}: combo, stile
      di gameplay (offensivo o stealth), equipaggiamento; di solito
      reversibili};
    \node[gd15sub, below=0.2cm of m1] (m2) {\textbf{Risorse}: pozioni,
      consumabili, punti statistica; spese permanenti o legate a spazi e
      momenti di gioco};
    \node[gd15sub, below=0.35cm of x] (x1) {\textbf{Tecnologia}: periferica,
      controller, arcade stick};
    \node[gd15sub, below=0.2cm of x1] (x2) {\textbf{Player behavior}: trolling,
      spam di frasi, comportamento verso altri umani};
    \node[gd15sub, below=0.2cm of x2] (x3) {\textbf{Community}: Reddit,
      Discord, side market (CS, Dota, workshop di Steam)};
  \end{tikzpicture}
  \caption{La tassonomia dei tipi di decisione presentata a lezione. Le
  categorie non sono mutualmente esclusive: la stessa decisione può stare
  in più rami.}
  \label{gd15:fig:tipi}
\end{figure}

\subsection{Le decisioni legate al tema}

\begin{description}
  \item[Emotion driven] Tutto ciò che è \textbf{illogico} e che scelgo
        ``a pelle'': Megaton, che Diego fa saltare in aria perché ``tanto è
        una nucleare in più''. Sono decisioni \textbf{viscerali}, guidate
        dalle emozioni, e il risultato sono decisioni che \textbf{il
        designer non può controllare}. Può spingere in alcune direzioni con
        tecniche manipolative (che si vedranno in futuro), ma su queste ha
        ancora meno controllo che sulle altre.
  \item[Narrative (e sociali)] Tutto ciò che è legato al raggiungere o
        meno un \textbf{finale} o una \textbf{linea narrativa} possibile.
        Rientrano qui anche le decisioni legate alla \textbf{socialità}: se
        voglio legare o no con altri personaggi, se voglio fare la
        \emph{romance} con Megera in \emph{Hades}\footnote{Nella trascrizione
        ``Medusa''. In \emph{Hades} c'è Dusa, una testa di gorgone con cui
        si può stringere un legame, mentre le romance vere e proprie sono con
        Megera e Thanatos; il docente intende in generale i rapporti
        sociali del gioco.} o con qualcun
        altro, e così via.
  \item[Estetiche] Qui ``estetica'' nel senso dell'\textbf{arte}, del
        \emph{dressing} (sezione~\ref{gd15:sec:target}). \textbf{Dove
        esploro}, se per predilezione personale preferisco il bioma innevato
        alla palude o al deserto; oppure una build \textbf{Fashion Souls}, la
        build di \emph{Dark Souls} che non punta a massimizzare le
        statistiche ma a essere \textbf{esteticamente appagante}.
\end{description}

\paragraph{Le build di Monster Hunter.} Nessuno degli studenti gioca a
\emph{Monster Hunter}, ma il docente mostra due build ``e qui devo ammettere
la mia'': una che per lui è ``terrificante'', bruttissima esteticamente,
probabilmente ottimizzata e con ``pure i gioielli giusti'', ma che ``non si
sposa bene''; e un'altra che lui sceglierebbe ``tutta la vita'': stesso
tema, stesso mostro, stessa rarità, stessa palette di colori, anche qualora
fosse meno ottimizzata. Non tutti i giocatori ragionano così; ma quando si
parla di Fashion Souls (il nome viene da \emph{Dark Souls}), le decisioni
legate alla moda, a come viene percepita l'armatura, ricadono qui, e sono
diverse dalle decisioni meccaniche.

\subsection{Le decisioni legate alle meccaniche}

Se invece voglio \textbf{massimizzare} la build, o adottare una strategia
specifica, o la \textbf{migliore strategia possibile}, siamo a cavallo tra
le meccaniche e, ancora di più, le decisioni meta. Le meccaniche si
dividono in due.
\begin{description}
  \item[Strategia] Che tipo di \textbf{combo} usare, il tipo di
        \textbf{gameplay} (più offensivo o più stealth, se il gioco lo
        permette), tutti gli \textbf{equipaggiamenti}: le strategie ad alto
        livello. Generalmente \textbf{possono sempre cambiare}: da momento a
        momento, da missione a missione, posso cambiare stile o
        equipaggiamento.
  \item[Risorse] Se spendo o no le risorse. Gran parte dei giocatori
        \textbf{tiene le pozioni}, i revitalizzanti, non spende i
        consumabili: è una decisione meccanica sulle risorse. Ci rientra
        anche come investo i \textbf{punti statistica}, come faccio crescere
        il personaggio. Qui la scelta, di solito, \textbf{non si annulla}:
        una volta bevuta, la pozione è consumata; una volta investito, il
        punto di norma non lo riottengo, a meno di tornare al falò o, come in
        \emph{Salt and Sanctuary}, di \textbf{resettare l'albero delle
        abilità}.
\end{description}
La differenza: una riguarda le \textbf{strategie} che il giocatore adotta,
l'altra le \textbf{risorse} che spende. E la spesa può essere
\textbf{permanente} oppure \textbf{temporanea}, ma di solito legata a spazi o
momenti di gioco: non si può cambiare l'esito della scelta immediatamente,
come invece si può fare con l'equipaggiamento.

\subsection{Le meta-decisioni}

L'ultimo tipo sono le \textbf{meta-decisioni}: tutto ciò che è legato alla
\textbf{tecnologia} fuori dal gioco (il controller, visto prima) e al
\textbf{comportamento del giocatore verso altri esseri umani} reali.
\begin{itemize}
  \item \textbf{Trollare} gli altri giocatori, cioè avere un obiettivo
        diverso da quello stabilito dal gioco, mettere i bastoni fra le
        ruote per interesse personale, umorismo, \emph{excitement}.
  \item Lo \textbf{spam} di una frase per annoiare, infastidire, frustrare
        gli altri: il meme di \emph{Overwatch}.\footnote{Nella trascrizione
        ``knee-dealing è meme di Overwatch''; il riferimento è con ogni
        probabilità alla ripetizione ossessiva di una linea vocale (per
        esempio ``I need healing''), un meme noto del gioco.}
  \item La \textbf{community}: tutto ciò che sta fuori dal gioco, Reddit,
        Discord, e il \textbf{side market}, la compravendita di oggetti al
        di fuori del gioco (\emph{Counter-Strike}, \emph{Dota}, gli oggetti
        del workshop di Steam).
\end{itemize}
È la componente sociale, il \emph{player behavior}: decisioni meta. Si
torna su tutte nel riepilogo finale (sezione~\ref{gd15:sec:leve}).

% ---------------------------------------------------------------------
\section{Le tre trappole: meaningless, obvious, blind}
\label{gd15:sec:trappole}

Torniamo alla slide delle porte: come si rende quella scelta una scelta che
\textbf{abbia senso implementare}, cioè che valga i costi e il tempo che le
si vogliono dedicare in produzione? Generalmente si seguono \textbf{tre
linee guida}, che sono altrettanti \textbf{must not} generali
(tabella~\ref{gd15:tab:trappole}).

\begin{table}[htbp]
  \centering
  \small
  \begin{tabular}{@{}p{0.16\textwidth}p{0.36\textwidth}p{0.4\textwidth}@{}}
    \toprule
    \textbf{Trappola} & \textbf{Che cosa succede} & \textbf{Esempi a lezione} \\
    \midrule
    \emph{Meaningless} & Una scelta va fatta, ma non ha un risultato
      effettivo nel gameplay, o lo ha molto più piccolo di quanto viene
      percepito. & Il ciondolo gabbia/uccello di \emph{BioShock Infinite};
      la colazione di \emph{Life is Strange}. \\
    \emph{Obvious} & Una delle opzioni è palesemente la migliore: la
      decisione non è una decisione (strategia dominante). & Il
      potenziamento della pistola in \emph{Control}; l'Anfora
      dell'Avidità in \emph{Yu-Gi-Oh!}. \\
    \emph{Blind} & Il giocatore sceglie senza informazioni sufficienti
      sull'effetto della scelta. & Le tre porte; i doni iniziali del primo
      \emph{Dark Souls}. \\
    \bottomrule
  \end{tabular}
  \caption{Le tre trappole delle decisioni. Sono linee guida: si può
  romperne una, non tutte, e solo dopo averne padroneggiato l'uso.}
  \label{gd15:tab:trappole}
\end{table}

\subsection{Decisioni meaningless}

La prima cosa da evitare (ed era il caso della porta) sono le decisioni
\textbf{senza senso}: c'è una scelta da fare, ma \textbf{non ha un effettivo
risultato nel gameplay}. Le porte, a dire il vero, un risultato ce l'hanno
avuto: uno \textbf{educativo} (un esempio da non seguire da cui partire) e
uno \textbf{umoristico}, ``che trova il tempo che trova''. Ma gli studenti non
hanno guadagnato nulla di materiale: i soldi e le tigri non c'erano. Nei
videogiochi vale lo stesso, ed è un discorso di \textbf{percezione}: se si
implementa una scelta percepita come \textbf{estremamente impattante}, ma
che nel gameplay non cambia nulla o pochissimo, probabilmente si sono
\textbf{investite male le risorse di produzione}.

\paragraph{Il ciondolo di BioShock Infinite.} Il docente chiede a Valerio se
ricordi la scelta del ciondolo per Elizabeth, \textbf{gabbia} o
\textbf{uccello}. Valerio la ricorda: ``mi pare che influiva sul finale''.
``Bella guess'', ma in realtà \textbf{non influiva su niente}; anche il
docente, la prima volta, l'aveva percepita così. È una scelta che la gente
ricorda (o ricorda male) come estremamente importante, ma che nel gioco è
\textbf{solo di \emph{flavor}}.

Per chi non l'ha giocato, il docente la ricostruisce (e cerca il video, ``che
ogni anno faccio fatica a beccare perché è sempre nascosto''): l'incipit di
\emph{BioShock Infinite} è andare in una città che non è più sott'acqua ma
\textbf{nel cielo}, per salvare una ragazza, Elizabeth, rinchiusa in una
torre con un carceriere enorme a forma di uccello (``il Big Daddy
volante'').\footnote{Il carceriere è Songbird.} C'è una simbologia molto
forte: il carceriere ha la forma di un uccello, e c'è il concetto di
\textbf{libertà}, perché lei ha vissuto tutta l'infanzia \textbf{in gabbia}.
Il giocatore deve scegliere per lei uno tra due ciondoli: quello con la
\textbf{gabbia} o quello con l'\textbf{uccello}. E la scelta viene presentata
in maniera molto importante:
\begin{itemize}
  \item \textbf{blocca il gameplay}: crea un \emph{funnel}, il level design
        si interrompe;
  \item c'è una \textbf{cutscene}, con personaggi nuovi, mai mostrati prima
        (che poi ritornano come nemici, ``se non sbaglio come boss'');
  \item è proprio una \textbf{presentazione} che il giocatore percepisce
        come importante.
\end{itemize}
Eppure da lì in avanti non serve più a nulla: cambia il ciondolo, ``che non
vedi mai'', perché giochi in prima persona e lei, per metà del gameplay, è
``lontana un chilometro''. È, molto probabilmente, un \textbf{refuso della
produzione}. La regola: \textbf{se il giocatore si pone davanti a una
scelta, rispetto a quanto importante viene percepita, deve avere un
risultato altrettanto importante nel gameplay}.

\subsection{Decisioni ovvie}

Il secondo punto sono le \textbf{decisioni ovvie}. A volte si può
suggerire al giocatore che una decisione è migliore di un'altra; ma avere
una \textbf{catena} di decisioni ovvie è sbagliato. Se la scelta ideale è
ovvia, il problema si pone; a maggior ragione se la scelta è
\textbf{permanente} e il giocatore sbaglia (per non aver trovato le
informazioni corrette o per qualsiasi altra ragione) e poi si rende conto che
una decisione ovvia era stata sbagliata: fa \emph{double down} sul proprio
senso di colpa. In \emph{Control}, per esempio, c'è un combattimento in cui un
potenziamento specifico di una pistola è \textbf{estremamente ovvio}, e con
quello si può ``cheesare'' lo scontro in maniera estrema.

\paragraph{L'Anfora dell'Avidità.} Ma l'esempio che il docente vuole citare,
``perché ha segnato la mia vita da grande amante di giochi di carte'', è
l'\textbf{Anfora dell'Avidità} (\emph{Pot of Greed}) di \emph{Yu-Gi-Oh!}.
Chi la conosce? (``Meno 10 punti a Grifondoro. Non è vero, sto
scherzando.'') Nella \textbf{banlist} è lì, tra le carte proibite, da
moltissimi anni (dopo un po' di ricerca a schermo, è Alberto a trovarla:
``non so come mai sono completamente cieco, grazie mille''). Il suo effetto è
semplicissimo: \textbf{peschi due carte}. Tu giochi l'Anfora e peschi due.
Questo fa.

Carte che fanno esattamente questo esistono anche in \emph{Magic} o nel GCC
di \emph{Pokémon}, eppure non sono così forti come in \emph{Yu-Gi-Oh!}, per
una serie di motivi:
\begin{itemize}
  \item il mazzo di \emph{Yu-Gi-Oh!} è \textbf{estremamente ridotto}
        rispetto a quello di \emph{Magic}: 40 carte;\footnote{Il docente
        dice ``40 rispetto a 80, minimo'', cioè ``la metà''. Il minimo per un
        mazzo constructed di \emph{Magic} è in realtà di 60 carte; il punto
        (un mazzo di \emph{Yu-Gi-Oh!} molto più piccolo, quindi più
        consistente) resta valido.}
  \item \emph{Yu-Gi-Oh!} è un gioco di \textbf{probabilità}, sempre più
        spiccatamente.
\end{itemize}
Metà degli studenti aveva analizzato proprio questo gioco
(``cosa sto dicendo? voi avete fatto l'analisi''), e il docente chiede: in
che direzione è spinto il gameplay oggi, ``ancora più di ieri''? Qual è il
problema di affrontare un avversario? Risponde Valerio: chi comincia
\textbf{inizia a pescare attivando effetti a catena} e controlla la partita,
con molte più possibilità di vincere. ``Assolutamente sì.'' I mazzi sono
ottimizzati per fare \textbf{più combo possibili con meno pezzi possibili}, e
per \textbf{non far giocare l'avversario}. Il docente scorre la banlist:
\begin{itemize}
  \item il \textbf{Maestro dello Shock}:\footnote{Nella trascrizione
        ``Maestro Shock''; è \emph{Shock Master}.} dichiari un tipo di carta,
        e quel tipo non può essere attivato; lo gioco nel mio turno, dichiaro
        ``mostro'', e l'avversario non gioca;
  \item \textbf{Bagooska}:\footnote{Nella trascrizione ``Bagusca'': è
        \emph{Number 41: Bagooska the Terribly Tired Tapir}.} finché è in
        posizione di difesa, gli effetti dei mostri non si possono attivare
        (``recentemente messo qua dentro'');
  \item \textbf{Rhongomyniad}:\footnote{Nella trascrizione ``Rongominion'' e
        ``Rongomidian'': è \emph{Number 86: Heroic Champion -
        Rhongomyniad}.} un mostro Xyz con cinque materiali; una volta per
        turno distrugge tutte le carte dell'avversario, è immune agli effetti
        delle carte, e quando ha quattro o più materiali (``e ce li ha'')
        l'avversario non può evocare normalmente o specialmente. Chi non
        gioca? L'avversario.
\end{itemize}
Gran parte delle carte proibite o creano combo o sono, da sole, carte che
\textbf{bloccano le giocate dell'avversario} (i cosiddetti
\emph{floodgate}).\footnote{Nella trascrizione ``floatgate''.} L'Anfora non
lo fa. Il problema è che oggi \textbf{metà del mazzo} è formata da
carte\footnote{Nella trascrizione ``Entrap''; dal contesto (carte che
bloccano l'avversario o abilitano le proprie giocate) sono quasi certamente
le \emph{hand trap}.} che servono a bloccare l'avversario o ad abilitare le
proprie giocate senza farlo giocare; e l'altra metà, a maggior ragione, deve
essere estremamente ottimizzata.

In questo contesto, avere l'Anfora significa che \textbf{ogni mazzo ce l'ha},
in \textbf{tre copie} (``manco in due''): è una carta così vantaggiosa (fa
pescare due carte, quindi aumenta di uno il numero di carte in mano) e si
gioca \textbf{senza costo}. Perché non dovrei giocarla? E perché non dovrei
metterla nel mazzo, se nello stato attuale del gioco \textbf{ogni carta è una
risorsa ovunque si trovi}: in mano, sul terreno, nel cimitero, a volte
persino nella zona bandita, perché posso recuperarla con un'altra carta o
attivarne l'effetto dal cimitero? In un gioco in cui la strategia è
\textbf{costruire il mazzo}, l'Anfora rende la scelta di inserirla non solo
ovvia, ma \textbf{necessaria}, in qualsiasi mazzo, aggro o control che sia.
E questo è il problema.

\subsection{Decisioni blind}

Il terzo punto sono le decisioni \textbf{blind}, ed è di nuovo il caso della
porta. Se il docente avesse detto ``la porta rossa ha il sacco di soldi'',
la decisione sarebbe stata \emph{ovvia}; non avendo dato nessuna
informazione, è il contrario: \textbf{scegliere senza sapere che effetto
avrà la scelta}. Non bisogna sempre esplicitare la ricompensa o dare tutte le
informazioni, ma bisogna \textbf{creare la condizione per cui il giocatore
sente il peso della scelta}, lo sente correttamente, o sente che si sta
impegnando per qualcosa per cui verrà premiato. È uno \textbf{spettro}: la
quantità di informazioni \textbf{sufficienti} da dare al giocatore per
metterlo in condizione di scegliere.

\paragraph{I doni iniziali di Dark Souls.} Un caso al tempo stesso negativo e
positivo, nel suo gioco: dopo la scelta della classe, il primo \emph{Dark
Souls} chiede di scegliere un oggetto da una lista, i \textbf{doni
iniziali}\footnote{Il docente dice che ``dovrebbero chiamarsi
\emph{starting gifts}'' e cerca la lista a schermo, mostrando prima per
errore quella di \emph{Dark Souls II}.} (tra cui la \textbf{Master Key} e
l'\textbf{Anello della Vecchia Strega}). Il giocatore non sa che cosa
facciano, o non ha informazioni sufficienti per capire se sia meglio la
chiave o l'anello rispetto alla classe scelta e alla strategia che vuole
adottare: \textbf{nel gioco non c'è ancora entrato}. La scelta diventa invece
estremamente interessante quando il giocatore \textbf{ha} informazioni sul
gioco, cioè alla seconda, terza, quarta run; e ha senso in un mondo come
quello di \emph{Dark Souls}, estremamente punitivo, cupo e spietato, che si
fregia di non dare tutte le soluzioni al giocatore: è parte della visione
del gioco. Applicata a \emph{Metroid} o ad altri titoli, la stessa scelta
sarebbe ``veramente problematica''. Dipende da dove la si fa.

\paragraph{Le regole si rompono, una alla volta.} Queste sono \textbf{must
not come linee guida generali}. Come sempre, se ne può \textbf{rompere una},
non tutte, e solo dopo averne \textbf{padroneggiato l'implementazione}.

\paragraph{Approfondimento: da dove vengono le tre trappole.} La terna
\emph{obvious}/\emph{meaningless}/\emph{blind} è quella del capitolo sugli
elementi di abilità strategica di \emph{Challenges for Game Designers}
(Brathwaite e Schreiber), che aggiunge qualche strumento utile.
\begin{itemize}
  \item Per le decisioni \textbf{ovvie} l'esempio è lo ``Highest Number
        Game'' (vince chi dice il numero più alto: il secondo giocatore non
        ha davvero una scelta). Il rimedio suggerito è spesso
        \textbf{rimuovere la scelta e automatizzarla} (i GDR moderni non
        chiedono più di mangiare e riposare, molti FPS ricaricano da soli),
        oppure aggiungere \textbf{pressione temporale}, trasformandola in una
        prova di destrezza.
  \item Per le decisioni \textbf{meaningless} il libro cita i vecchi GDR:
        ``Vuoi vedere il re?'' Se rispondi no, ``Ma devi!'', e lo vedi
        comunque. Con una precisazione che anticipa la discussione della
        sezione successiva: la \textbf{percezione}. Molti giochi moderni
        offrono scelte narrative che non cambiano l'esito, ma la prima
        volta il giocatore percepisce che lo fanno, e \textbf{quella
        percezione rende la scelta significativa}; ci si accorge del
        contrario solo rigiocando.
  \item Per le decisioni \textbf{blind} l'esempio è la roulette, che resta
        avvincente solo perché c'è in palio denaro vero. Il rimedio è dare
        \textbf{abbastanza} informazioni, non tutte: basta una voce
        sentita prima su un NPC per trasformare una scelta alla cieca in una
        scelta interessante.
\end{itemize}
Schell, in \emph{The Art of Game Design}, chiama \textbf{strategia
dominante} il caso in cui un'opzione è chiaramente migliore delle altre
(``una volta scoperta, il gioco non è più divertente, perché il puzzle è
risolto''), e ne fa la domanda chiave della \textbf{Lente n.~39, delle
scelte significative}: quali scelte chiedo di fare al giocatore? Sono
significative, e come? Gliene do il numero giusto? Ci sono strategie
dominanti? Sul numero, Schell riporta l'osservazione di Michael Mateas:
se le scelte superano i desideri del giocatore, questi si sente
\textbf{sopraffatto}; se sono meno, \textbf{frustrato}; se coincidono, prova
\textbf{libertà e appagamento}, che è lo stesso equilibrio delle tre porte
(e dei tre gusti) del capitolo~\ref{gd9:cap:agency-visione}.

% ---------------------------------------------------------------------
\section{La discussione: e la ricompensa emotiva?}
\label{gd15:sec:discussione}

\subsection{L'obiezione di Federica}

``Dubbi? Chiaro?'' Una smorfia di Federica non convince il docente. Lei non
conosce \emph{BioShock}, ma chiede: una scelta senza conseguenze reali non
rischia comunque di influire sulla \textbf{filosofia} del gioco, come qui la
gabbia? Se scegli la gabbia, magari lei fa una smorfia e poi si continua.

In effetti, risponde il docente, c'è un \textbf{bark narrativo} di Elizabeth
(gli studenti non l'hanno sentito, c'erano solo i sottotitoli): una frase che
risponde alla scelta dell'uno o dell'altro pendente. Può essere più una
\textbf{ricompensa emotiva} che concreta, e se ne è discusso anche in altri
anni. Il docente sarebbe d'accordo, se non fosse per la \textbf{bilancia}:
da una parte la \textbf{messa in scena} (tutto ciò che porta alla scelta e i
costi che ci sono dietro: la cutscene, i personaggi nuovi, la nuova camera),
dall'altra il \textbf{risultato}, cioè un bark su gabbia o uccello. La messa
in scena è percepita come così importante rispetto al risultato che la
decisione, a conti fatti, \textbf{viene percepita come meaningless}. È questa
la discriminante.

Federica insiste con un esempio inventato: un momento di festa, la ragazza
appena liberata, un regalo a lei: la ricompensa emotiva potrebbe valere la
messa in scena. Il docente è d'accordo, e propone una soluzione da designer:
se si volesse mantenere quella scelta ma farla percepire come \textbf{meno
importante}, la si potrebbe mettere in una \textbf{parte opzionale}. Non
blocca il gameplay né il flusso; chi va dritto (o fa speedrun) la può
perdere; chi guarda i dettagli (ed è il primo momento in cui si vede
davvero la struttura sociale della città nel cielo, con la spiaggia
simulata, le onde d'acqua finta, l'estetica dell'America degli anni
Venti) viene premiato trovandola non sulla strada principale, ma in un
angolo (non nascosto, né misterioso). Ciò che non convince il docente è la
messa in scena finale. E, ripete, probabilmente non è nemmeno voluta: è un
\textbf{evento tagliato}, finito nella build così perché non c'è stato tempo
di produzione per espanderlo altrove; ``i costi sono andati in quella
direzione, quello è stato''.

Ci possono essere scelte meaningless più azzeccate come esempio (``mi sono
dimenticato di aggiornare le slide; lo farò''): il docente propone di
\textbf{trovarne una} e discuterne venerdì.

\subsection{Il canto per la maestra: Silksong}

Per mostrare che il risvolto emotivo può essere una soluzione, il docente
chiede a Francesca, che ha finito \emph{Silksong}, se abbia visto la
\textbf{morte della maestra di Shakra},\footnote{Nella trascrizione
``chakra''; Shakra è la cartografa di \emph{Silksong}, ``quella che ti vende
le mappe''.} cioè la sub-quest legata alla venditrice di mappe. È una
sub-quest estremamente lunga: presuppone che il giocatore abbia ottenuto la
mappa di tutte le zone, e quindi abbia trovato Shakra in tutte le zone. E
alla fine la ricompensa (al netto di un'arma, ``se non erro'' gli anelli) è
\textbf{una canzone}: sei tu, con Shakra, a fare l'ultimo canto alla maestra
morta. La ricompensa è principalmente \textbf{emotiva}, non meccanica; eppure
il docente, come giocatore, si è sentito appagato: era la fine di un percorso
in cui un'amica (``immedesimandomi in Hornet'') cercava la sua maestra nel
labirinto e poi la trovava morta. ``Estremamente struggente.'' Avere la
certezza della sua morte e poterle dare l'ultimo saluto è stato
sufficiente: una risoluzione \textbf{emotiva}, non meccanica.

\subsection{Valerio: ``ho chiamato un amico''}

Valerio rinforza il punto su \emph{BioShock}: quando ci stava giocando si è
fermato per contattare l'amico che gli aveva consigliato la serie e
chiedergli se quella scelta avesse una grande influenza. La messa in scena
funzionava, quindi: era proprio un discorso di \textbf{percezione} del
giocatore. Il docente è comunque dell'idea che esistano esempi ancora più
costosi e ancora più meaningless, e rimanda a venerdì.

\subsection{Francesca: e la colazione?}

Francesca si collega al discorso emotivo: seguendo questa linea di pensiero,
anche \textbf{pancake o uova e bacon} in \emph{Life is Strange} non è una
decisione meaningless? ``Sono perfettamente d'accordo'': il docente l'aveva
portata come esempio di \emph{decisione}, ma non ha mai detto che fosse
meaningful. Francesca immagina che serva a dare al giocatore la percezione di
essere libero di scegliere, cioè \textbf{agency}. Esatto, e c'è di più:
\begin{itemize}
  \item \emph{Life is Strange} è un gioco \textbf{estremamente
        cinematografico}, in cui per gran parte dell'esperienza il giocatore
        guarda cutscene; prima di quella scelta ci saranno stati almeno
        \textbf{sette minuti} di cutscene (lei che si sveglia, va a svegliare
        l'amica, vanno alla tavola calda con la madre, parlano). C'è
        l'immedesimazione, e c'è la \textbf{rottura del pace}: dopo sette
        minuti senza scelte il giocatore ``si addormenta'', per quanto
        narrativo sia il gioco e per quanto adatta sia l'audience;
  \item la variante è \textbf{minima a livello di produzione}: la cutscene
        resta la stessa, le animazioni sono le stesse, cambia solo
        \textbf{l'asset nel piatto};
  \item il giocatore la vede come \textbf{estremamente minore} rispetto a
        salvare una persona. In \emph{Life is Strange} ci sono momenti in
        cui le scelte decidono la vita o la morte di altri NPC, e la
        differenza si percepisce nettamente.
\end{itemize}
Il docente confessa di essere ``uno di quei malati di mente'' che nei giochi
narrativi cerca di min-maxare, prendendo le decisioni eticamente più
corrette, soprattutto al primo playthrough: faceva dieci minuti di gioco,
vedeva che cosa succedeva e ricaricava; e il gioco ha pure la meccanica del
\textbf{tornare indietro nel tempo} (``ci vado a nozze''). Come giocatore, gli
era sempre chiaro quando stava facendo una scelta minore e quando una
maggiore. Una scelta, insomma, dipende \textbf{dalla presentazione} ma anche
da \textbf{quanto costa metterla a terra}.

\subsection{Il caso limite: ``Clementine will remember that''}

C'è un caso limite, ``che non vi dovrei nemmeno portare perché ho paura che
vi confonda'', dello stile \textbf{Telltale}, maestra dei giochi narrativi a
episodi: \emph{The Walking Dead}. Il docente mostra prima una scena di
\emph{Life is Strange: Before the Storm} (``opinione impopolare: a me è
piaciuto di più lo spin-off''), con la notifica a schermo che la scelta avrà
\textbf{conseguenze}. In \emph{The Walking Dead} questa meccanica è ancora
più spinta, al punto da essere diventata un meme: \textbf{``Clementine will
remember that''}. E in abbastanza casi la scelta \textbf{non è davvero
impattante} e non cambia niente nel gameplay.

Perché funziona? Con \textbf{un elemento di UI} si fa percepire un'emozione
estremamente potente. In mezzo all'apocalisse zombie il protagonista si porta
dietro una bambina, e le sue scelte la \textbf{educano}: salvare o non
salvare una persona, scegliere tra la vita e la morte, raccontarle bugie su
chi è morto. Come designer, si vuole che il giocatore senta il \textbf{peso
della responsabilità} nella relazione tra il protagonista e l'NPC di punta:
deve sempre sentirsi osservato, sentire che le sue azioni hanno conseguenze
su una minorenne. E lo si vuole \textbf{con meno costi possibili}: lo stesso
gioco con tutte le scelte davvero meaningful si potrebbe fare, ma sarebbe
estremamente difficile da produrre. È una \textbf{scelta furba}: con poco il
giocatore ottiene tanto, ``anche troppo'', tanto che molti (il docente
compreso) rifacevano lo stesso capitolo più volte per capire se cambiava
qualcosa. Lo stesso vale in \emph{Life is Strange}: alcune scelte segnalate
come ricordate non verranno ricordate affatto; quelle davvero impattanti
hanno la notifica, ma la notifica non implica una scelta impattante.

È un \textbf{caso limite di rottura} di queste linee guida, fatto da
un'azienda che faceva \emph{solo} quello (episodi spiccatamente narrativi con
quel sistema) ed era quindi estremamente verticale su
quello.\footnote{Telltale Games ha chiuso nel 2018. Il docente aggiunge che
chi ci lavorava ``è in Adobe Studio'': quasi certamente AdHoc Studio, fondato
da ex membri di Telltale; il nome nella trascrizione non è affidabile.}

A questo punto Daniele deve andare a prendere il treno e saluta (``buon
weekend a tutti'').

% ---------------------------------------------------------------------
\section{Il contrario: decisioni ben costruite}
\label{gd15:sec:buone}

Il contrario delle tre trappole, ovviamente, è ciò che va fatto
(tabella~\ref{gd15:tab:buone}).

\subsection{Non meaningless}

Scelte che modificano, anche di poco, l'\textbf{allineamento} di altri
personaggi o portano a un finale diverso. L'esempio, già visto in passato,
è \textbf{\emph{Oxenfree}}:\footnote{Nella trascrizione ``Aux Free''.}
ogni scelta modifica leggermente l'allineamento dei personaggi e porta a uno
o più finali. Se con tuo fratello\footnote{In \emph{Oxenfree} è il
fratellastro, Jonas.} sei stato molto (o anche poco più che poco)
maleducato, il legame si rompe; due NPC possono mettersi insieme o no a
seconda di certe azioni, ma \textbf{mancarne una o due} e fare tutte le
altre porta comunque a quel risultato. E ci sono finali specifici
\emph{ad hoc}: far tornare il fratello morto\footnote{Michael, il fratello
della protagonista Alex.} è possibile solo se si è rotto un certo tipo di
\textbf{loop}, che è alla base di questo ``teen drama, teen horror,
coming-of-age horror, come volete chiamarlo''. Lo stesso vale per
\textbf{\emph{The Witcher}}: se tratti Ciri\footnote{Nella trascrizione
``Searing''; dal contesto (``se le trattate in una certa maniera, poi
compirà il sacrificio finale'') è quasi certamente Ciri in \emph{The Witcher
3}, il cui destino dipende da come Geralt si è comportato con lei.} in una
certa maniera, poi compirà il sacrificio finale, e così via. C'è una scelta,
e ha un impatto \textbf{allineato con l'aspettativa del giocatore} e con come
il designer la mette in scena.

\subsection{Non ovvie}

C'è una decisione da fare, ma \textbf{non c'è una strategia dominante}. Non
a caso gli esempi sono spiccatamente strategici e lato build: la possibilità
di modificare le strategie a seconda del proprio desiderio da giocatore.

\paragraph{Fire Emblem.} Chi ha giocato un \emph{Fire Emblem}? Il docente
consiglia di recuperarli (``non sarò certo io a dirvi di scaricare emulatori
per GBA e far girare la ROM di \emph{Blazing Blade}'', ma se lo si facesse
\ldots). Si ha un \textbf{party} che cresce con la trama; i personaggi si
muovono su \textbf{griglia}; ci sono obiettivi di missione in missione;
alcuni nemici si possono \textbf{reclutare}; i personaggi hanno un'anima da
JRPG, con design e personalità spiccate; nei capitoli recenti si può fare
\emph{romance}, o farla fare ai personaggi, in base a chi si porta in
battaglia. I personaggi hanno statistiche e possono \textbf{cambiare classe}
(a schermo c'è \emph{The Sacred Stones}): cambiando classe aumenta una
statistica e cambia la crescita, come l'evoluzione dei Pokémon. E il
giocatore \textbf{non porta necessariamente i più forti}, né li fa evolvere
necessariamente nella classe più forte: magari vuole quel personaggio in
squadra per la romance, o vuole vedere le relazioni sociali tra due
personaggi, o (nei Fire Emblem più vecchi) solo tra un lord e il suo
cavaliere di punta.

\paragraph{Pokémon.} Snorlax ha statistiche che puntano chiaramente a un
Pokémon \textbf{difensivo}, un muro fisico; eppure lo si può buildare per
spingerlo in \textbf{velocità}, per renderlo imprevedibile, o dargli quattro
mosse inattese dall'avversario, a seconda anche delle meccaniche della
generazione (Megaevoluzione, Gigamax e così via). Una \textbf{via insolita}
rispetto allo stile di gioco: questo significa decisioni non ovvie.

\subsection{Non blind}

Decisioni che hanno un effetto sul gioco, e il giocatore ha
\textbf{sufficiente conoscenza} per capire quale. Di nuovo \emph{Life is
Strange}: il \textbf{finale} è scegliere chi salvare. Anche se non viene
detto esplicitamente, aver giocato fino a quel momento ha dato informazioni
sufficienti su che cosa comporta la scelta, perché è stata
\textbf{costruita in tutte le ore precedenti} (``assumo dieci ore di
gameplay''). Questo la rende non blind.

\begin{table}[htbp]
  \centering
  \small
  \begin{tabular}{@{}p{0.2\textwidth}p{0.36\textwidth}p{0.36\textwidth}@{}}
    \toprule
    \textbf{Obiettivo} & \textbf{Come} & \textbf{Esempi} \\
    \midrule
    Non meaningless & Impatto allineato con l'aspettativa e con la messa
      in scena; anche piccoli spostamenti di allineamento contano. &
      \emph{Oxenfree}, \emph{The Witcher 3} (Ciri). \\
    Non ovvie & Nessuna strategia dominante; più vie valide a seconda del
      desiderio del giocatore. & \emph{Fire Emblem} (classi, squadre,
      relazioni), \emph{Pokémon} (build insolite). \\
    Non blind & Informazioni sufficienti, anche implicite, costruite nel
      tempo. & Il finale di \emph{Life is Strange}. \\
    \bottomrule
  \end{tabular}
  \caption{Le decisioni ben costruite, contrario delle tre trappole.}
  \label{gd15:tab:buone}
\end{table}

% ---------------------------------------------------------------------
\section{Il riepilogo: le leve di ogni tipo di decisione}
\label{gd15:sec:leve}

Il docente si prende ``giusto un minuto'' per scusarsi (il flowchart non ci
starà) e ricapitola. Le categorie principali sono tre, non necessariamente
mutualmente esclusive; per ognuna ci sono le \textbf{leve} su cui il
designer può agire (tabella~\ref{gd15:tab:leve}).

\begin{table}[htbp]
  \centering
  \small
  \begin{tabular}{@{}p{0.15\textwidth}p{0.38\textwidth}p{0.39\textwidth}@{}}
    \toprule
    \textbf{Tipo} & \textbf{Leva} & \textbf{In pratica} \\
    \midrule
    Tema & \emph{Awareness} & Il giocatore sa che cosa sta scegliendo e
      quanto pesa. \\
         & Conseguenze & Tempo e risorse investiti (una romance, un
      personaggio trattato male) devono avere una conseguenza. \\
         & \emph{Reminder} & La notifica a scelta fatta (Clementine): più è
      alto il reminder, più il giocatore deve percepirne il peso. \\
         & Permanenza vs \emph{trade-off} & Uccidere una Sorellina (e il
      finale) contro l'armatura che oggi indosso e domani posso cambiare. \\
    \midrule
    Meccaniche & Scelta delle azioni & Tante azioni possibili, una per
      volta (la lotta Pokémon). \\
         & Rischio vs ricompensa & Alto rischio, alta ricompensa; premiare
      i prudenti e far sperare chi cerca adrenalina. \\
         & Lungo vs breve termine & Spendere ora o risparmiare per dopo (i
      soldi di \emph{Hollow Knight}). \\
         & Scambio di risorse & Oro per una pozione: il ciclo economico. \\
    \midrule
    Meta & Bluff, \emph{dealmaking}, \emph{backstabbing} & Giocare onesti o
      disonesti con altri umani (\emph{Among Us}, \emph{Sea of Thieves}). \\
         & \emph{Min-maxing} & Conoscere il gioco oltre quanto il designer
      si aspetta (i rendimenti decrescenti dei Souls). \\
    \bottomrule
  \end{tabular}
  \caption{Il riepilogo finale della lezione: le leve delle decisioni per
  tipo.}
  \label{gd15:tab:leve}
\end{table}

\subsection{Tema}

Il tema si porta dietro le decisioni \textbf{emozionali, viscerali e di
istinto}, e le leve sono quattro.
\begin{description}
  \item[Awareness] Il giocatore deve sapere \textbf{che cosa sta
        scegliendo} e \textbf{quanto peso ha} la scelta. Il ciondolo di
        \emph{BioShock Infinite} si può categorizzare come decisione lato
        tema, solo che la percezione è molto più alta del risultato
        effettivo.
  \item[Conseguenze] Le scelte devono avere \textbf{conseguenze}: se investo
        tempo e risorse per una romance, mi aspetto una conseguenza di
        qualche tipo; se prendo tutte le scelte giuste, o se tratto
        terribilmente un personaggio in \emph{Fallout} e ho una statistica
        di ``cattiveria'' altissima (``non mi ricordo come si
        chiama''),\footnote{È il \emph{Karma} di \emph{Fallout 3}.} la
        relazione con gli altri personaggi ne risente.
  \item[Reminder] Il caso di Clementine: i pop-up a schermo. Più è alto il
        reminder, più il giocatore deve percepire la pesantezza e l'effetto
        della scelta. E il reminder va dato \textbf{quando la scelta è stata
        fatta}: hai scelto, te lo ricordo, è importante.
  \item[Permanenza vs trade-off] Una scelta \textbf{permanente} (uccidere una
        Sorellina, che influenza il finale) contro un \textbf{trade-off}
        (scelgo uno o l'altro). Di nuovo Fashion Souls: posso avere addosso
        una sola armatura, ma oggi ho questa e domani posso cambiarla.
\end{description}

\subsection{Meccaniche}

\begin{description}
  \item[Scelta delle azioni] È alla base di ogni tipo di gameplay, più o
        meno esteso: al giocatore si mettono in mano una serie di azioni, ma
        \textbf{ne può fare una per volta}. In una lotta Pokémon ho quattro
        mosse, oppure la borsa e la Poké Ball, oppure la fuga; non posso
        farle tutte, e farne una esclude l'altra: se fuggo, quel Pokémon
        l'ho perso; se sconfiggo lo \emph{shiny} che ho davanti, non lo posso
        più catturare. Questo rende una scelta più importante, o percepita
        come tale.
  \item[Rischio vs ricompensa] Il classico: alto rischio e alta ricompensa,
        basso rischio e bassa ricompensa. Si possono premiare i giocatori più
        prudenti e al tempo stesso far sperare quelli in cerca di emozioni
        adrenaliniche.
  \item[Lungo vs breve termine] Investo adesso e perdo la risorsa in futuro?
        Come con Federica: in \emph{Hollow Knight} non spendo i soldi adesso
        perché punto a qualcosa di più costoso, oppure li spendo tutti per
        non rischiare di bruciarmeli morendo contro un boss.
  \item[Scambio di risorse] Pago oro per una pozione, spendo qualcosa per
        ottenere qualcos'altro: il \textbf{ciclo economico} lato meccanico.
        L'ultima lezione dell'anno sarà dedicata proprio
        all'\textbf{economia}.
\end{description}

\paragraph{Approfondimento: la triangolarità.} Schell dedica al rischio
contro ricompensa una lente intera, la \textbf{Lente n.~40, della
triangolarità}: dare al giocatore la scelta tra giocare sul sicuro per una
ricompensa bassa o rischiare per una alta. L'esempio è \emph{Space
Invaders}, con il disco volante che vale molti punti ma costringe a esporsi.
\emph{Challenges for Game Designers} osserva inoltre che chi è in svantaggio
tende a rischiare di più per recuperare, e chi è in vantaggio a giocare sul
sicuro per conservarlo: una dinamica che genera tensione da sola. Lo stesso
libro distingue il \textbf{trade-off} (non ci sono risorse per tutti gli
obiettivi: arma migliore o armatura migliore?) dal \textbf{dilemma}, in cui
tutte le opzioni danneggiano il giocatore.

\subsection{Meta}

Tutto ciò che, ripete il docente, è legato alla \textbf{socialità} e alla
relazione con altri individui reali (e sulla slide c'è apposta l'omino di
\emph{Among Us}): \textbf{bluffing}, \textbf{dealmaking},
\textbf{backstabbing}. Si può scegliere di adottarli, che il gioco sia di
deduzione sociale o no: giocare onestamente o disonestamente, per esempio
nel dividere il bottino; in \emph{Sea of Thieves} ci si può permettere di
essere onesti o disonesti con gli altri giocatori.

Qui rientra anche il \textbf{min-maxing}, l'estrema conseguenza della scelta
delle azioni: conoscere le meccaniche al punto da \textbf{ottimizzare
qualsiasi guadagno e minimizzare qualsiasi spesa}. Perché sta tra le meta e
non tra le meccaniche? Perché presuppone una conoscenza del gioco
\textbf{molto più profonda di quella che il designer si aspetta}. Il docente
lascia la ricerca agli studenti: gli \emph{i-frame} in \emph{Monster Hunter}
o in \emph{Dark Souls}, il min-max nei Souls, la spesa delle anime in
\emph{Dark Souls}. Si scopre, con dei grafici, che da un certo punto in
avanti diventa \textbf{sconveniente investire punti} in una statistica,
perché l'aumento non è più vantaggioso come prima
(figura~\ref{gd15:fig:softcap}); vale per quasi tutti i soulslike di
FromSoftware (``non tutti i soulslike''). Ma è un calcolo che il designer non
si aspetta che il giocatore faccia; e i giocatori di questo tipo stanno, sulla
curva normale (sezione~\ref{gd10:sec:gaussiana}), in una parte molto spinta
verso gli \textbf{evangelisti}.

\begin{figure}[htbp]
  \centering
  \begin{tikzpicture}[x=0.11cm, y=0.042cm]
    \draw[-{Stealth[length=2mm]}] (0,0) -- (62,0);
    \node[font=\footnotesize] at (31,-17) {punti investiti nella statistica};
    \draw[-{Stealth[length=2mm]}] (0,0) -- (0,110) node[above,
      font=\footnotesize] {valore effettivo};
    \draw[very thick, blue!60!black] (0,0) -- (20,70) -- (40,92) -- (58,98);
    \draw[dashed, gray] (20,0) -- (20,70);
    \draw[dashed, gray] (40,0) -- (40,92);
    \node[below, font=\scriptsize] at (20,0) {soglia 1};
    \node[below, font=\scriptsize] at (40,0) {soglia 2};
    \node[font=\scriptsize, align=center] at (5,72) {molto\\conveniente};
    \node[font=\scriptsize, align=center] at (30,100) {meno\\conveniente};
    \node[font=\scriptsize, align=center] at (51,108) {quasi\\inutile};
  \end{tikzpicture}
  \caption{Lo schema dei rendimenti decrescenti (\emph{soft cap}) di cui
  parla il docente: oltre certe soglie, ogni punto investito rende sempre
  meno. Grafico illustrativo con valori inventati, non i dati di un gioco
  specifico.}
  \label{gd15:fig:softcap}
\end{figure}

% ---------------------------------------------------------------------
\section{Per la prossima lezione}
\label{gd15:sec:prossima}

La lezione successiva riprenderà da qui, agganciandosi
all'\textbf{audience}: si è già visto un primo strumento, se ne vedranno
altri, e anche il principio \emph{data informed, non data driven} (già
sfiorato parlando di feedback, sezione~\ref{gd10:sec:gaussiana}). Il
\textbf{flowchart di Grotte} si farà la settimana successiva, se c'è tempo;
altrimenti il docente si rimette alle domande degli studenti. ``Mi spiace non
averlo fatto, colpa mia.'' Resta comunque un buon punto di partenza, meno
sfidante di \emph{Pac-Man}, per testare le conoscenze sui flowchart; e se
gli studenti vogliono giocarci in una serata, il docente lo consiglia.
``Teniamola, per favore, la porta aperta'': per qualsiasi domanda sui
flowchart, scrivergli.

% ---------------------------------------------------------------------
\section{Il task: le decisioni di un gioco ``truffa''}
\label{gd15:sec:task}

\subsection{Il consiglio: un gioco di venti minuti}

``So che voi state bramando un task, e io vi do il miglior task che potrei
mai darvi''. Prima, un consiglio: un gioco dello stesso designer di
\emph{The Stanley Parable},\footnote{Il docente lo chiama ``The Emerald and
the Tiger, qualcosa'': è quasi certamente \emph{Dr. Langeskov, The Tiger,
and The Terribly Cursed Emerald: A Whirlwind Heist} (Crows Crows Crows,
2015), gratuito su Steam e lungo circa un quarto d'ora, di William Pugh,
coautore di \emph{The Stanley Parable}.} gratuito su Steam, che dura circa
venti minuti: da giocare per il discorso sulla \textbf{percezione}, anche in
vista del secondo anno. Ma non è quello il gioco del task.

\subsection{La consegna}

Il gioco del task è ``ancora migliore'', dello stesso autore:\footnote{Il
titolo non compare nella trascrizione: il docente lo manda nella chat di
Zoom. Dalla descrizione è un gioco comico fatto con RPG Maker, lungo una
decina di minuti, in cui Keanu Reeves viene accusato di un crimine e deve
scagionarsi; non è stato possibile identificarlo con certezza. Il docente
avverte anche che all'apertura l'audio è assordante.} un gioco
\textbf{comico}, fatto con \textbf{RPG Maker}, che dura una
\textbf{decina di minuti}, in cui \textbf{Keanu Reeves} viene accusato
(ingiustamente o giustamente, a seconda) di un crimine e deve scagionarsi.
Attenzione all'apertura: l'audio è assordante. La consegna:
\begin{enumerate}
  \item giocare il gioco;
  \item fare la \textbf{lista delle decisioni} presenti nel gioco,
        \textbf{categorizzate per tipo} (tema, meccaniche, meta, con le
        sottocategorie): a ogni decisione il suo tipo, o i suoi tipi;
  \item \textbf{motivare}, da designer, perché ogni decisione \textbf{è o
        non è} \emph{meaningless}, \emph{obvious} e \emph{blind}.
\end{enumerate}

\subsection{Le mani avanti}

``È una truffa, questa. È un'estrema truffa.'' In questo gioco umoristico
molte decisioni che normalmente ricadrebbero in questi tipi, e che
normalmente sarebbero meaningful, sono \textbf{completamente sovvertite}:
per esempio (``giusto per darvi questo indizio'') la \textbf{scelta del
personaggio}. Bisogna quindi prestare attenzione: se lo si osserva con
l'occhio giusto, si vede che molto di ciò che normalmente ricade in una
categoria può essere \textbf{ribaltato}, a seconda del gioco e del designer.

\subsection{La chiusura}

``Vi lascio la scoperta di questa incredibile esperienza.'' Il docente
ripete l'invito a vedersi in presenza venerdì, per chi può, e la sua
disponibilità a parlare dentro e fuori dalla lezione; resta collegato per
eventuali domande, altrimenti lascia gli studenti ``liberi come Keanu Reeves
alla fine dell'esperienza''.

% ---------------------------------------------------------------------
\section*{Punti chiave del capitolo}
\addcontentsline{toc}{section}{Punti chiave del capitolo}

\begin{itemize}
  \item Niente più membri esterni in aula, salvo richiesta dell'intera
        classe; la registrazione è uno strumento per gli studenti; da remoto,
        webcam accese. Venerdì prossimo è una lezione importante: meglio in
        presenza.
  \item Posizioni di nemici e oggetti: non documenti separati. Il level
        design documenta a monte (level design document, mappa 2D, bubble
        map), ma il dettaglio (percorsi, distanze) sta \textbf{in engine},
        dopo il blockout, e c'è sempre una discrepanza tra 2D e 3D. Il
        funzionamento di un oggetto (il vaulting) sta nel game design; le
        quantità per zona stanno nel \textbf{Master Asset Document}, con
        hero asset e dipartimenti coinvolti.
  \item In \emph{Dread} lo \textbf{stealth non è un pilastro} (il Phantom
        Cloak perde valore, e le zone EMMI chiedono di uscire, non di
        nascondersi); \textbf{gli EMMI sì}: un predatore inattaccabile
        confinato in zone che fanno da spina dorsale ai livelli. Lo stealth è
        una \textbf{strategia}, non una meccanica; il combattimento è un
        \textbf{sistema}. USP: gli EMMI, e l'accorpamento ``capitolo
        conclusivo, 2D non Prime, su Switch''.
  \item Il target audience era un extra: i tipi di giocatori arrivano
        venerdì, e il gioco va mirato a uno o due archetipi. Non mescolare i
        due sensi di ``estetiche'' (aspetto, come Schell, o emozioni).
  \item Must do, must not e nice to have, lato design, sono l'elenco dei
        \textbf{sistemi e delle meccaniche} necessari, esclusi o accessori:
        camera, abilità, counter, EMMI e loro limiti, boss, mappa (must) e
        puntine (nice), biomi $\pm 1$, hard mode, collezionabili; la valuta
        spendibile è un must not di Metroid, l'\textbf{economy sink} un nice
        to have di \emph{Hollow Knight}. Non ripetere nei must not il
        contrario di un must do.
  \item Non capire non è un problema; non chiedere sì. La ricerca su giochi
        sconosciuti è voluta: guardare dieci secondi e annotare le
        meccaniche; dividere la qualità dell'attenzione.
  \item Si possono dividere i must per dipartimento; l'ambientazione
        (sci-fi) è \textbf{vestizione}. Test: leggere la lista e immaginare
        il gioco; quanto si avvicina all'originale? Camera (tre Zelda),
        personaggio, abilità (il loro numero dice i costi), struttura del
        mondo. Lo skill tree è must in \emph{Borderlands}, nice altrove.
  \item Gli esempi non sono la regola: sperimentare. Stesso mandate,
        feedback diversi: la risposta giusta la dice solo la produzione.
  \item Flowchart: niente blocchi slegati, niente rettangoli con più uscite
        (si usa il rombo). Il flowchart di Grotte è rimandato.
  \item Calendario: lezione venerdì 27; il 14 pitching con Diego e Jack e un
        task di gruppo da \textbf{presentare in aula il 18}, con feedback a
        voce. Il 27 novembre c'è il Beer to Peer a Milano.
  \item Le tre porte: tre scelte, ma nessuna informazione prima e nessuna
        sulla ricompensa. \textbf{Tutto è una decisione}, anche il
        controller (meta-decisione); in analisi interessano quelle che si
        possono controllare.
  \item Tipi di decisione, non esclusivi: \textbf{tema} (emotion driven,
        narrative e sociali, estetiche), \textbf{meccaniche} (strategia,
        reversibile; risorse, spesa permanente o temporanea),
        \textbf{meta} (tecnologia, player behavior, community e side
        market).
  \item Le tre trappole: \textbf{meaningless} (il ciondolo di \emph{BioShock
        Infinite}: messa in scena sproporzionata al risultato),
        \textbf{obvious} (l'Anfora dell'Avidità, necessaria in ogni mazzo),
        \textbf{blind} (le porte, i doni iniziali di \emph{Dark Souls}). Se ne
        rompe una, non tutte, e solo con padronanza.
  \item La ricompensa emotiva può bastare (la canzone di \emph{Silksong}), se
        la bilancia tra messa in scena e risultato regge; la colazione di
        \emph{Life is Strange} è meaningless ma economica e rompe il pace;
        ``Clementine will remember that'' è un caso limite di rottura,
        furbo e poco costoso.
  \item Il contrario: non meaningless (\emph{Oxenfree}, \emph{The Witcher
        3}), non ovvie (\emph{Fire Emblem}, Pokémon), non blind (il finale di
        \emph{Life is Strange}).
  \item Le leve: tema (awareness, conseguenze, reminder, permanenza vs
        trade-off), meccaniche (scelta delle azioni, rischio vs ricompensa,
        lungo vs breve termine, scambio di risorse), meta (bluff,
        dealmaking, backstabbing, min-maxing).
  \item Task: giocare il gioco comico con Keanu Reeves, elencare e
        categorizzare le decisioni e motivare se sono meaningless, obvious,
        blind; attenzione, è una ``truffa'' che sovverte le aspettative.
\end{itemize}
